\documentclass[12pt,a4paper]{article}
\usepackage[T1]{fontenc}
\usepackage[utf8]{inputenc}
\usepackage{lmodern}
\usepackage{amsmath,amssymb,amsthm,mathtools}
\usepackage{geometry}
\geometry{margin=1in}
\usepackage{microtype}
\usepackage{hyperref}
\usepackage{xcolor}
\usepackage{enumitem}
\usepackage{fancyhdr}
\usepackage{bookmark}

\usepackage{geometry}
\geometry{margin=1in}
\usepackage{microtype}
\usepackage{enumitem}
\usepackage{hyperref}
\usepackage{array}
\usepackage{booktabs}
\usepackage{xcolor}

\usepackage{booktabs}
\usepackage{longtable}
\usepackage{enumitem}
\usepackage{hyperref}

\hypersetup{
	colorlinks=true,
	linkcolor=blue,
	urlcolor=blue,
	citecolor=blue
}

\hypersetup{
	colorlinks=true,
	linkcolor=blue!50!black,
	urlcolor=blue!50!black,
	citecolor=blue!50!black,
	pdftitle={Theory of Commutative Algebra 7},
	pdfauthor={},
	pdfsubject={Solved problems and theory notes in commutative algebra},
	pdfcreator={OpenAI}
}

\setlist[itemize]{topsep=4pt,itemsep=2pt,parsep=0pt,partopsep=0pt}
\setlength{\parskip}{0.55em}
\setlength{\parindent}{0pt}

\setcounter{secnumdepth}{2}


\setcounter{tocdepth}{2}

\pagestyle{fancy}


\fancyhf{}
\lhead{Theory of Commutative Algebra 12}
\rhead{\thepage}
\renewcommand{\headrulewidth}{0.4pt}

\newtheoremstyle{notespace}{6pt}{6pt}{\normalfont}{} {\bfseries}{.}{0.5em}{}
\theoremstyle{notespace}
\newtheorem{definition}{Definition}[section]
\newtheorem{remark}{Remark}[section]


\newcommand{\Spec}{\operatorname{Spec}}
\newcommand{\Frac}{\operatorname{Frac}}
\newcommand{\OO}{\mathcal O}
\newcommand{\kk}{\mathbf{k}}
\newcommand{\A}{\mathbb A}
\newcommand{\Z}{\mathbb Z}
\newcommand{\Q}{\mathbb Q}
\newcommand{\R}{\mathbb R}
\newcommand{\C}{\mathbb C}
\newcommand{\F}{\mathbb F}
\newcommand{\mf}{\mathfrak}
\newcommand{\m}{\mathfrak m}
\newcommand{\p}{\mathfrak p}

\theoremstyle{definition}

\newtheorem{example}[definition]{Example}



\DeclareMathOperator{\Ass}{Ass}
\DeclareMathOperator{\Ann}{Ann}
\DeclareMathOperator{\Supp}{Supp}
\DeclareMathOperator{\Hom}{Hom}
\DeclareMathOperator{\Min}{Min}
\DeclareMathOperator{\Ker}{Ker}


\DeclareMathOperator{\depth}{depth}
\DeclareMathOperator{\Tor}{Tor}
\DeclareMathOperator{\Ext}{Ext}

\DeclareMathOperator{\coker}{coker}
\DeclareMathOperator{\im}{im}
\newcommand{\Coker}{\operatorname{Coker}}



\newtheorem{proposition}[definition]{Proposition}
\newtheorem{theorem}[definition]{Theorem}
\newtheorem{corollary}[definition]{Corollary}


\newcommand{\kerf}{\operatorname{Ker}}

\newcommand{\pd}{\operatorname{pd}}
\newcommand{\rank}{\operatorname{rank}}
\newcommand{\syz}{\operatorname{Syz}}
\newcommand{\ZZ}{\mathbb{Z}}







\hypersetup{
	colorlinks=true,
	linkcolor=blue!60!black,
	urlcolor=blue!60!black,
	citecolor=blue!60!black
}

\title{\vspace{-1.5em}\bfseries Theory of Commutative Algebra 12}
\author{}
\date{\today}

\begin{document}
	\maketitle
	\thispagestyle{plain}
		\begin{center}
		\begin{minipage}{0.92\textwidth}
			Lecture notes with worked examples in commutative algebra
		\end{minipage}
	\end{center}
		\pdfbookmark[1]{Contents}{contents}\tableofcontents
	\clearpage
\section{Residue Fields of Points of a Scheme}

One of the most important local invariants of a point on a scheme is its \emph{residue field}.  
It tells us what the ``scalars at the point'' are, and it is one of the basic tools for understanding:

\begin{itemize}[leftmargin=2em]
	\item rational points,
	\item generic points,
	\item fibers of morphisms,
	\item tangent spaces,
	\item function fields of irreducible varieties.
\end{itemize}

The point of the construction is that a point of a scheme is not merely a geometric location in the naive sense.  
A point corresponds to a prime ideal, and the residue field measures what remains of the local ring after one mods out by the maximal ideal of that local ring.

In affine geometry this is completely explicit, and in general it glues well across affine charts.

\subsection{Definition of the residue field}

\subsubsection{Affine case}

Let
\[
X=\Spec(A)
\]
be an affine scheme, and let
\[
x\in X
\]
correspond to a prime ideal
\[
\p\subseteq A.
\]

Then the local ring of \(X\) at \(x\) is
\[
\OO_{X,x}=A_{\p}.
\]
Its maximal ideal is
\[
\m_x=\p A_{\p}.
\]

\begin{definition}
	The \emph{residue field} of the point \(x\) is
	\[
	k(x):=\OO_{X,x}/\m_x=A_{\p}/\p A_{\p}.
	\]
\end{definition}

So one first localizes at the prime \(\p\), obtaining the local ring \(A_{\p}\), and then quotients by its maximal ideal.

\subsubsection{General scheme case}

Now let \(X\) be an arbitrary scheme and let \(x\in X\).  
Choose an affine open neighborhood
\[
U=\Spec(A)\subseteq X
\]
containing \(x\). Then \(x\) corresponds to a prime ideal
\[
\p\subseteq A,
\]
and one defines
\[
\OO_{X,x}=A_{\p},\qquad k(x)=A_{\p}/\p A_{\p}.
\]

This does not depend on the chosen affine neighborhood \(U\), because the local ring \(\OO_{X,x}\) is intrinsic to the scheme \(X\).

\subsection{Why \(k(x)\cong \Frac(A/\p)\)}

The most useful computational formula is the following.

\begin{proposition}
	Let \(A\) be a ring and let \(\p\subseteq A\) be a prime ideal.  
	Then
	\[
	A_{\p}/\p A_{\p}\cong \Frac(A/\p).
	\]
	Hence, for the point \(x\in \Spec(A)\) corresponding to \(\p\),
	\[
	k(x)\cong \Frac(A/\p).
	\]
\end{proposition}

\begin{proof}
	Consider the exact sequence
	\[
	0\longrightarrow \p \longrightarrow A \longrightarrow A/\p \longrightarrow 0.
	\]
	Localizing at the multiplicative set
	\[
	S=A\setminus \p
	\]
	gives an exact sequence
	\[
	0\longrightarrow S^{-1}\p \longrightarrow S^{-1}A \longrightarrow S^{-1}(A/\p)\longrightarrow 0.
	\]
	But
	\[
	S^{-1}A=A_{\p}
	\]
	and
	\[
	S^{-1}\p=\p A_{\p}.
	\]
	So
	\[
	A_{\p}/\p A_{\p}\cong S^{-1}(A/\p).
	\]
	
	Now \(A/\p\) is an integral domain, because \(\p\) is prime.  
	The image of \(S=A\setminus \p\) in \(A/\p\) is exactly the set of nonzero elements of \(A/\p\).  
	Therefore localizing \(A/\p\) at that set produces its field of fractions:
	\[
	S^{-1}(A/\p)\cong \Frac(A/\p).
	\]
	Hence
	\[
	A_{\p}/\p A_{\p}\cong \Frac(A/\p).
	\]
\end{proof}

\begin{remark}
	This formula is the quickest way to compute residue fields in practice:
	\[
	\boxed{
		k(x)=\Frac(A/\p).
	}
	\]
	
	So the computation usually has two steps:
	\begin{enumerate}[leftmargin=2em]
		\item compute the quotient \(A/\p\),
		\item take its fraction field.
	\end{enumerate}
\end{remark}

\subsection{Geometric meaning}

\subsubsection{What does the residue field measure?}

The residue field \(k(x)\) is the field in which functions can be evaluated at the point \(x\).

At an ordinary \(k\)-rational point of a variety over a field \(k\), one gets
\[
k(x)=k.
\]
At a non-rational closed point, one often gets a finite extension of \(k\).  
At a generic point, one gets a function field.

Thus the residue field distinguishes three very different kinds of behavior:

\begin{itemize}[leftmargin=2em]
	\item \textbf{closed rational points:} \(k(x)=k\),
	\item \textbf{closed non-rational points:} \(k(x)\) is an algebraic extension of \(k\),
	\item \textbf{generic points:} \(k(x)\) is a rational function field.
\end{itemize}

\subsubsection{\(k\)-rational points}

\begin{definition}
	Let \(X\) be a scheme over a field \(k\).  
	A point \(x\in X\) is called \emph{\(k\)-rational} if
	\[
	k(x)\cong k.
	\]
\end{definition}

For example, on \(\A^1_k=\Spec(k[t])\), the point corresponding to the maximal ideal \((t-a)\) with \(a\in k\) is \(k\)-rational, because its residue field is \(k\).

\subsubsection{Generic points}

If \(X\) is irreducible and \(\eta\) is its generic point, then \(\eta\) corresponds to the zero ideal in an affine integral coordinate ring.  
Its residue field is the field of rational functions on \(X\).

So for an integral affine scheme \(X=\Spec(A)\), where \(A\) is a domain, the generic point has residue field
\[
k(\eta)=\Frac(A).
\]

This is why residue fields are so central in algebraic geometry.

\subsection{Basic examples}

\subsubsection{Example 1: affine line over a field}

Let
\[
X=\Spec(k[t]),
\]
where \(k\) is a field.

Prime ideals in \(k[t]\) are of two basic types:

\begin{itemize}[leftmargin=2em]
	\item the zero ideal \((0)\),
	\item ideals generated by irreducible polynomials.
\end{itemize}

\paragraph{Closed rational point}

Take the prime ideal
\[
\p=(t-a),\qquad a\in k.
\]
Then
\[
k[t]/(t-a)\cong k,
\]
so
\[
k(x)=\Frac(k[t]/(t-a))\cong \Frac(k)\cong k.
\]

Thus the point \(t=a\) is a \(k\)-rational point.

\paragraph{Generic point}

Take
\[
\p=(0).
\]
Then
\[
k(x)=\Frac(k[t]/(0))=\Frac(k[t])=k(t).
\]

So the generic point of the affine line has residue field \(k(t)\), the rational function field in one variable.

\paragraph{Closed non-rational point}

If \(k\) is not algebraically closed, there may be maximal ideals generated by irreducible polynomials of degree \(>1\).  
For example, over \(\R\), the ideal \((t^2+1)\) is maximal, and its residue field is \(\C\); this is treated in detail below.

\subsubsection{Example 2: affine line over \(\R\)}

Let
\[
X=\Spec(\R[t]).
\]

\paragraph{Real rational point}

Take
\[
\p=(t-a),\qquad a\in \R.
\]
Then
\[
\R[t]/(t-a)\cong \R,
\]
hence
\[
k(x)=\R.
\]

So these are the ordinary real points.

\paragraph{A closed point that is not \(\R\)-rational}

Now take
\[
\p=(t^2+1).
\]
Since \(t^2+1\) is irreducible in \(\R[t]\), this ideal is prime, in fact maximal. Then
\[
\R[t]/(t^2+1)\cong \C.
\]
Since this quotient is already a field, its fraction field is itself:
\[
k(x)=\Frac(\R[t]/(t^2+1))\cong \C.
\]

This is a very important example.  
It shows that a closed point of a scheme over \(\R\) does \emph{not} have to have residue field \(\R\).  
It may have residue field a finite algebraic extension of \(\R\).

Geometrically, the polynomial \(t^2+1\) has no real root, but scheme-theoretically it still defines a closed point of \(\Spec(\R[t])\).

\subsubsection{Example 3: \(\Spec(\Z)\)}

Let
\[
X=\Spec(\Z).
\]

The prime ideals of \(\Z\) are:
\[
(0),\ (2),\ (3),\ (5),\dots
\]

\paragraph{Closed point \((p)\)}

Take a prime number \(p\). The corresponding point is the prime ideal
\[
\p=(p).
\]
Then
\[
\Z/(p)\cong \F_p,
\]
so
\[
k(x)=\Frac(\Z/(p))=\F_p.
\]

Thus every closed point of \(\Spec(\Z)\) has residue field a finite field.

\paragraph{Generic point}

Take the prime ideal
\[
\p=(0).
\]
Then
\[
k(x)=\Frac(\Z)=\Q.
\]

So \(\Spec(\Z)\) has a single generic point with residue field \(\Q\), and all the closed points have residue fields \(\F_p\).

This is one of the basic pictures in scheme theory:
\[
\text{generic point } \leftrightarrow \Q,\qquad
\text{closed point } (p)\leftrightarrow \F_p.
\]

\subsubsection{Example 4: affine plane over a field}

Let
\[
X=\Spec(k[s,t]).
\]

\paragraph{A rational closed point}

Take
\[
\p=(s-a,t-b),\qquad a,b\in k.
\]
Then
\[
k[s,t]/(s-a,t-b)\cong k,
\]
and hence
\[
k(x)=k.
\]

So an ordinary point \((a,b)\in k^2\) corresponds to a maximal ideal whose residue field is \(k\).

\paragraph{Generic point of an irreducible curve}

Take
\[
\p=(t-s^2).
\]
Then
\[
k[s,t]/(t-s^2)\cong k[s],
\]
so
\[
k(x)=\Frac(k[s])=k(s).
\]

This point is not a closed point.  
It is the generic point of the irreducible parabola
\[
t=s^2.
\]

So the generic point of an irreducible subvariety carries its function field.

\paragraph{Generic point of the whole plane}

Take
\[
\p=(0).
\]
Then
\[
k(x)=\Frac(k[s,t])=k(s,t).
\]

So the generic point of the whole affine plane has residue field the rational function field in two variables.

\subsubsection{Example 5: a reducible scheme \(k[x,y]/(xy)\)}

Let
\[
A=k[x,y]/(xy),\qquad X=\Spec(A).
\]
Geometrically, this is the union of the two coordinate axes in the plane.

There are several important points on this scheme.

\paragraph{Generic point of the \(y\)-axis}

Consider the prime ideal
\[
\p=(x)\subseteq A.
\]
Then
\[
A/\p \cong \frac{k[x,y]/(xy)}{(x)} \cong k[y].
\]
Therefore
\[
k(x)=\Frac(k[y])=k(y).
\]

So \((x)\) is the generic point of the \(y\)-axis.

\paragraph{Generic point of the \(x\)-axis}

Similarly, if
\[
\p=(y),
\]
then
\[
A/\p \cong k[x],
\qquad
k(x)=k(x).
\]

Thus \((y)\) is the generic point of the \(x\)-axis.

\paragraph{The intersection point}

Take the maximal ideal
\[
\m=(x,y)\subseteq A.
\]
Then
\[
A/\m \cong k,
\]
so
\[
k(x)=k.
\]

Thus the crossing point is a \(k\)-rational closed point.

\begin{remark}
	This example is extremely useful because it shows three different kinds of points on the same scheme:
	\begin{itemize}[leftmargin=2em]
		\item generic point of one irreducible component,
		\item generic point of the other irreducible component,
		\item a closed intersection point.
	\end{itemize}
	Their residue fields are correspondingly:
	\[
	k(y),\qquad k(x),\qquad k.
	\]
\end{remark}

\subsubsection{Example 6: affine line over a finite field}

Let
\[
X=\Spec(\F_q[t]).
\]

A closed point corresponds to a maximal ideal generated by an irreducible polynomial
\[
f(t)\in \F_q[t].
\]
Then
\[
k(x)=\F_q[t]/(f(t)).
\]

If \(\deg(f)=d\), then
\[
\F_q[t]/(f(t))\cong \F_{q^d}.
\]
Thus
\[
k(x)\cong \F_{q^d}.
\]

So closed points of \(\Spec(\F_q[t])\) do not correspond only to elements of \(\F_q\).  
They also correspond to irreducible polynomials of higher degree, and their residue fields are finite extensions of \(\F_q\).

This is a standard and very important point in arithmetic geometry.

\subsubsection{Example 7: local ring viewpoint at a rational point}

Let
\[
A=k[t],\qquad \p=(t-a).
\]
Then
\[
A_{\p}
\]
consists of fractions
\[
\frac{f(t)}{g(t)}
\]
with
\[
g\notin (t-a),
\]
equivalently,
\[
g(a)\neq 0.
\]

Its maximal ideal is
\[
\p A_{\p}=(t-a)A_{\p}.
\]
Then
\[
A_{\p}/(t-a)A_{\p}\cong k.
\]

This is a concrete way to think of the residue field: functions that are regular at \(a\), modulo those that vanish at \(a\), leave only the value at \(a\).

\subsection{Examples with detailed computations}

\subsubsection{Detailed computation in \(\Spec(k[t])\)}

Take
\[
A=k[t],\qquad \p=(t-a).
\]
We compute:
\[
A/\p \cong k[t]/(t-a)\cong k.
\]
Hence
\[
k(x)=\Frac(A/\p)=\Frac(k)=k.
\]

If instead \(\p=(0)\), then
\[
A/\p=k[t],
\]
so
\[
k(x)=\Frac(k[t])=k(t).
\]

Thus in one and the same scheme \(\Spec(k[t])\) we see:
\[
\text{closed rational points } \rightsquigarrow k,\qquad
\text{generic point } \rightsquigarrow k(t).
\]

\subsubsection{Detailed computation in \(\Spec(\R[t])\) at \((t^2+1)\)}

Take
\[
A=\R[t],\qquad \p=(t^2+1).
\]
Then
\[
A/\p = \R[t]/(t^2+1).
\]
Since \(t^2+1\) is irreducible in \(\R[t]\), the quotient is a field.  
In fact,
\[
\R[t]/(t^2+1)\cong \C
\]
via the class of \(t\mapsto i\).

Therefore
\[
k(x)=\Frac(A/\p)=A/\p\cong \C.
\]

So this point is closed but not \(\R\)-rational.

\subsubsection{Detailed computation in \(\Spec(\Z)\)}

Take
\[
A=\Z,\qquad \p=(p)
\]
for a prime number \(p\). Then
\[
A/\p=\Z/(p)=\F_p.
\]
Since \(\F_p\) is already a field,
\[
k(x)=\F_p.
\]

If instead \(\p=(0)\), then
\[
A/\p=\Z,
\]
and
\[
k(x)=\Frac(\Z)=\Q.
\]

\subsubsection{Detailed computation for the generic point of a curve}

Let
\[
A=k[s,t],\qquad \p=(t-s^2).
\]
Then
\[
A/\p \cong k[s,t]/(t-s^2)\cong k[s].
\]
Since \(k[s]\) is a domain, its fraction field is
\[
\Frac(k[s])=k(s).
\]
Hence
\[
k(x)=k(s).
\]

Thus the generic point of the parabola carries the function field of the parabola.

\subsection{Rational points versus closed points}

In elementary algebraic geometry over an algebraically closed field \(k\), people often identify points with tuples of coordinates in \(k^n\).  
Scheme theory generalizes this.

\begin{itemize}[leftmargin=2em]
	\item A closed point of \(\Spec(A)\) corresponds to a maximal ideal.
	\item A \(k\)-rational point is a closed point whose residue field is exactly \(k\).
	\item A closed point need not be \(k\)-rational if the base field is not algebraically closed.
\end{itemize}

For example, in \(\Spec(\R[t])\):
\[
(t-a)\quad \text{has residue field } \R,
\]
but
\[
(t^2+1)\quad \text{has residue field } \C.
\]

So every \(k\)-rational point is a closed point, but not every closed point is \(k\)-rational.

\subsection{Residue fields and generic points of irreducible subsets}

Let \(X\) be a scheme and let \(Y\subseteq X\) be an irreducible closed subset.  
Then \(Y\) has a unique generic point \(\eta_Y\).

If \(X=\Spec(A)\) is affine and \(Y=V(\p)\) for a prime ideal \(\p\), then the generic point of \(Y\) is precisely the point corresponding to \(\p\), and
\[
k(\eta_Y)=\Frac(A/\p).
\]

So the residue field of the generic point of an irreducible closed subset is exactly the function field of that subset.

This viewpoint is central in algebraic geometry.  
For example:

\begin{itemize}[leftmargin=2em]
	\item the generic point of \(\A^1_k\) has residue field \(k(t)\),
	\item the generic point of the parabola \(t=s^2\) has residue field \(k(s)\),
	\item the generic point of the \(y\)-axis in \(V(xy)\) has residue field \(k(y)\).
\end{itemize}

\subsection{Why residue fields appear in many constructions}

\subsubsection{Fibers of morphisms}

If
\[
f:X\to Y
\]
is a morphism of schemes and \(y\in Y\), then the fiber over \(y\) is naturally a scheme over the residue field \(k(y)\), not merely over the original base field.

More precisely, if \(f\) is affine and corresponds to a ring map
\[
B\to A,
\]
and $(y)$ corresponds to a prime ideal $(q \subseteq B)$, then the fiber is
\[
X_y \cong \Spec\!\bigl(A\otimes_B k(y)\bigr).
\]

So the residue field of the base point tells us over which field the fiber should be considered.



\subsubsection{Tangent spaces}

The Zariski tangent space at a point \(x\in X\) is naturally a vector space over \(k(x)\):
\[
T_xX \cong \operatorname{Hom}_{k(x)}\!\bigl(\m_x/\m_x^2,\;k(x)\bigr).
\]

Thus even tangent vectors live over the residue field of the point, not automatically over the original ground field.

This is another reason why residue fields are fundamental rather than secondary.

\subsection{A summary table of standard examples}

\begin{center}
	\renewcommand{\arraystretch}{1.35}
	\begin{longtable}{>{\raggedright\arraybackslash}p{0.24\textwidth}
			>{\raggedright\arraybackslash}p{0.24\textwidth}
			>{\raggedright\arraybackslash}p{0.19\textwidth}
			>{\raggedright\arraybackslash}p{0.23\textwidth}}
		\toprule
		\textbf{Scheme \(X\)} & \textbf{Point / prime ideal} & \textbf{Type of point} & \textbf{Residue field \(k(x)\)} \\
		\midrule
		\endfirsthead
		
		\toprule
		\textbf{Scheme \(X\)} & \textbf{Point / prime ideal} & \textbf{Type of point} & \textbf{Residue field \(k(x)\)} \\
		\midrule
		\endhead
		
		\bottomrule
		\endfoot
		
		\(\Spec(k[t])\) & \((t-a)\) & closed \(k\)-rational & \(k\) \\
		\(\Spec(k[t])\) & \((0)\) & generic point & \(k(t)\) \\
		\(\Spec(\R[t])\) & \((t-a)\) & closed \(\R\)-rational & \(\R\) \\
		\(\Spec(\R[t])\) & \((t^2+1)\) & closed non-rational & \(\C\) \\
		\(\Spec(\Z)\) & \((p)\) & closed point & \(\F_p\) \\
		\(\Spec(\Z)\) & \((0)\) & generic point & \(\Q\) \\
		\(\Spec(k[s,t])\) & \((s-a,t-b)\) & closed \(k\)-rational & \(k\) \\
		\(\Spec(k[s,t])\) & \((t-s^2)\) & generic point of parabola & \(k(s)\) \\
		\(\Spec(k[s,t])\) & \((0)\) & generic point of plane & \(k(s,t)\) \\
		\(\Spec(k[x,y]/(xy))\) & \((x)\) & generic point of \(y\)-axis & \(k(y)\) \\
		\(\Spec(k[x,y]/(xy))\) & \((y)\) & generic point of \(x\)-axis & \(k(x)\) \\
		\(\Spec(k[x,y]/(xy))\) & \((x,y)\) & closed intersection point & \(k\) \\
		\(\Spec(\F_q[t])\) & \((f(t))\), \(f\) irreducible, \(\deg f=d\) & closed point & \(\F_{q^d}\) \\
	\end{longtable}
\end{center}

\subsection{A compact conceptual summary}

For a point \(x\in X\), the residue field \(k(x)\) is the field obtained by taking the local ring at \(x\) and modding out by its maximal ideal:
\[
k(x)=\OO_{X,x}/\m_x.
\]

In the affine case \(X=\Spec(A)\), if \(x\leftrightarrow \p\), then
\[
k(x)=A_{\p}/\p A_{\p}\cong \Frac(A/\p).
\]

This single formula explains all the main cases:

\begin{itemize}[leftmargin=2em]
	\item ordinary rational points give the base field,
	\item closed non-rational points give field extensions,
	\item generic points give function fields.
\end{itemize}

So residue fields are the correct algebraic language for saying what it means to ``evaluate functions at a scheme-theoretic point.''

\subsection{Conclusion}

The residue field of a point is one of the simplest and most fundamental local invariants in scheme theory.

It is important because:

\begin{itemize}[leftmargin=2em]
	\item it is easy to compute:
	\[
	k(x)=\Frac(A/\p),
	\]
	\item it distinguishes rational points, non-rational closed points, and generic points,
	\item it appears naturally in fibers, tangent spaces, and function fields,
	\item it gives a precise algebraic meaning to the idea of values of functions at a point.
\end{itemize}

For many students, the first real conceptual shift in scheme theory is realizing that a point is not just a coordinate tuple; it comes equipped with a local ring and a residue field.  
Once this is understood, many later constructions become much more natural.


\subsection{Why \(k[s,t]/(s-a,t-b)\cong k\) and not \(k\times k\)}

A natural question is the following:

\medskip

Why do we have
\[
k[s,t]/(s-a,t-b)\cong k
\]
rather than
\[
k[s,t]/(s-a,t-b)\cong k\times k?
\]

\medskip

The reason is that the ideal
\[
(s-a,t-b)
\]
corresponds to \emph{one single point}, namely the point
\[
(a,b)\in k^2,
\]
so quotienting by this ideal means that both variables are forced to take those fixed values:
\[
s=a,\qquad t=b.
\]
After this quotient, every polynomial becomes just one scalar in \(k\), namely its value at the point \((a,b)\).

\subsubsection{The evaluation map}

Consider the ring homomorphism
\[
\varphi:k[s,t]\longrightarrow k
\]
defined by
\[
\varphi(f(s,t))=f(a,b).
\]
This is the evaluation map at the point \((a,b)\).

We claim that:
\[
\ker(\varphi)=(s-a,t-b).
\]

Indeed, a polynomial \(f(s,t)\) lies in the kernel if and only if
\[
f(a,b)=0.
\]
But this is equivalent to saying that \(f\) vanishes at the point \((a,b)\), and by the basic algebra of polynomial rings, this happens exactly when \(f\) belongs to the ideal generated by \(s-a\) and \(t-b\):
\[
f\in (s-a,t-b).
\]

Moreover, the map \(\varphi\) is surjective, because every element \(c\in k\) is the image of the constant polynomial \(c\).

Therefore, by the First Isomorphism Theorem,
\[
k[s,t]/(s-a,t-b)\cong k.
\]

\subsubsection{Direct substitution viewpoint}

Another way to see this is to interpret the quotient directly.

In the quotient ring
\[
k[s,t]/(s-a,t-b),
\]
the classes of \(s-a\) and \(t-b\) are zero. Hence
\[
s=a,\qquad t=b
\]
inside the quotient.

Therefore, every polynomial \(f(s,t)\) becomes equal to the scalar
\[
f(a,b)\in k.
\]

For example:
\[
s^2+3st+t \;\mapsto\; a^2+3ab+b.
\]
So the quotient ring contains no independent variables anymore. It consists only of scalars from \(k\). That is why it is isomorphic to \(k\).

\subsubsection{Why it is not \(k\times k\)}

The product ring
\[
k\times k
\]
appears when one is simultaneously keeping track of \emph{two separate points} or \emph{two separate components}. It does not correspond to one point.

To understand the difference, compare the following two situations.

\paragraph{One point}

The ideal
\[
(s-a,t-b)
\]
defines exactly one point \((a,b)\). Hence
\[
k[s,t]/(s-a,t-b)\cong k.
\]

\paragraph{Two points}

Now suppose one wants to keep track of two distinct points on the affine line, say \(t=a\) and \(t=b\) with \(a\neq b\). Then one considers
\[
k[t]/\bigl((t-a)(t-b)\bigr).
\]
Since the ideals \((t-a)\) and \((t-b)\) are comaximal, the Chinese Remainder Theorem gives
\[
k[t]/\bigl((t-a)(t-b)\bigr)
\cong
k[t]/(t-a)\times k[t]/(t-b)
\cong
k\times k.
\]

Here we really do have two distinct points, and the quotient remembers the value of a polynomial at both of them:
\[
f(t)\mapsto (f(a),f(b)).
\]

So the ring \(k\times k\) corresponds to two separate pieces of data, while \(k[s,t]/(s-a,t-b)\) corresponds to only one point.

\subsubsection{Field argument}

There is also a very quick abstract argument.

The ideal
\[
(s-a,t-b)
\]
is a maximal ideal in \(k[s,t]\). Indeed,
\[
k[s,t]/(s-a,t-b)\cong k,
\]
and \(k\) is a field, so the ideal is maximal.

Now if \(\mathfrak m\) is a maximal ideal in any ring \(A\), then
\[
A/\mathfrak m
\]
must be a field.

But the ring
\[
k\times k
\]
is \emph{not} a field, because for example
\[
(1,0)\neq 0
\]
but \((1,0)\) is not invertible.

Therefore
\[
k[s,t]/(s-a,t-b)
\]
cannot possibly be isomorphic to \(k\times k\), since the quotient by a maximal ideal must be a field.

\subsubsection{A concrete example}

Take \(k=\mathbb{R}\), \(a=1\), \(b=2\). Then
\[
\mathbb{R}[s,t]/(s-1,t-2)
\]
is the quotient in which
\[
s=1,\qquad t=2.
\]
So for example:
\[
s^2+t \mapsto 1^2+2=3,
\]
\[
st+5 \mapsto 1\cdot 2+5=7,
\]
\[
s^3-4t \mapsto 1-8=-7.
\]

Thus every polynomial becomes a real number. The quotient ring is exactly \(\mathbb{R}\), not \(\mathbb{R}\times \mathbb{R}\).

\subsubsection{Comparison with a quotient that really gives a product}

Let us contrast this with
\[
\mathbb{R}[t]/(t^2-1).
\]
Since
\[
t^2-1=(t-1)(t+1)
\]
and the ideals \((t-1)\) and \((t+1)\) are comaximal, the Chinese Remainder Theorem gives
\[
\mathbb{R}[t]/(t^2-1)\cong \mathbb{R}[t]/(t-1)\times \mathbb{R}[t]/(t+1)
\cong \mathbb{R}\times\mathbb{R}.
\]

This quotient does not correspond to one point. It corresponds to two reduced points:
\[
t=1\qquad \text{and} \qquad t=-1.
\]

So here the product ring appears because the quotient is keeping track of two points at once.

\subsubsection{Geometric interpretation}

Geometrically:

\begin{itemize}
	\item
	\[
	(s-a,t-b)
	\]
	is the ideal of the single point \((a,b)\subset \mathbb{A}^2_k\), so its coordinate ring is
	\[
	k.
	\]
	
	\item
	A ring like
	\[
	k\times k
	\]
	is the coordinate ring of two disjoint reduced points.
\end{itemize}

Thus:
\[
\text{one point} \longleftrightarrow k,
\qquad
\text{two separate points} \longleftrightarrow k\times k.
\]

\subsubsection{Final conclusion}

So the correct statement is:
\[
k[s,t]/(s-a,t-b)\cong k,
\]
because quotienting by \((s-a,t-b)\) means evaluating polynomials at the single point \((a,b)\).

It is \emph{not} isomorphic to \(k\times k\), because \(k\times k\) would correspond to two separate points, and also because \(k\times k\) is not a field, while the quotient by a maximal ideal must always be a field.

\medskip

\textbf{Solution.}
The quotient
\[
k[s,t]/(s-a,t-b)
\]
is obtained by imposing the relations
\[
s=a,\qquad t=b.
\]
Hence every polynomial \(f(s,t)\) becomes the scalar \(f(a,b)\in k\). Equivalently, the evaluation homomorphism
\[
\varphi:k[s,t]\to k,\qquad f(s,t)\mapsto f(a,b)
\]
is surjective and has kernel
\[
\ker(\varphi)=(s-a,t-b).
\]
Therefore, by the First Isomorphism Theorem,
\[
k[s,t]/(s-a,t-b)\cong k.
\]

By contrast, a product ring such as \(k\times k\) appears when one quotients by an ideal corresponding to two separate points, for example
\[
k[t]/((t-a)(t-b))\cong k\times k
\qquad (a\neq b),
\]
via the Chinese Remainder Theorem. Thus \(k[s,t]/(s-a,t-b)\) is \(k\), not \(k\times k\).

\subsection{Why \((t-a,s-b)\) is a single point in \(k[s,t]\)}

Let
\[
A=k[s,t].
\]
The indeterminates of this polynomial ring are
\[
s \quad \text{and} \quad t.
\]
They are the coordinate functions on the affine plane
\[
\mathbb{A}^2_k.
\]

Thus a point of \(\mathbb{A}^2_k\) in these coordinates is written as
\[
(s,t)=(\alpha,\beta)
\]
for some \(\alpha,\beta\in k\).

Now consider the ideal
\[
(t-a,s-b)\subseteq k[s,t].
\]
This means that in the quotient ring we impose the relations
\[
t=a,\qquad s=b.
\]
So the corresponding point is exactly
\[
(s,t)=(b,a).
\]

Hence \((t-a,s-b)\) is the ideal of the single point whose \(s\)-coordinate is \(b\) and whose \(t\)-coordinate is \(a\).

\subsubsection{Order of generators does not matter}

We may also write
\[
(t-a,s-b)=(s-b,t-a),
\]
because the ideal generated by two elements does not depend on the order in which the generators are written.

So both notations define the same ideal and hence the same point.

\subsubsection{Why this is one point}

The quotient
\[
k[s,t]/(t-a,s-b)
\]
forces both variables to become fixed scalars:
\[
s=b,\qquad t=a.
\]
Therefore every polynomial \(f(s,t)\) becomes the scalar
\[
f(b,a)\in k.
\]
There is no freedom left in either variable, so the quotient corresponds to exactly one point.

Equivalently, the evaluation homomorphism
\[
\varphi:k[s,t]\longrightarrow k,
\qquad
f(s,t)\longmapsto f(b,a)
\]
has kernel
\[
\ker(\varphi)=(s-b,t-a)=(t-a,s-b).
\]
Hence by the First Isomorphism Theorem,
\[
k[s,t]/(t-a,s-b)\cong k.
\]

Since the quotient is \(k\), this is the coordinate ring of a single reduced point.

\subsubsection{General rule}

More generally, in the polynomial ring
\[
k[x_1,\dots,x_n],
\]
the ideal
\[
(x_1-a_1,\dots,x_n-a_n)
\]
is the ideal of the single point
\[
(a_1,\dots,a_n)\in \mathbb{A}^n_k.
\]

Thus in the ring
\[
k[s,t],
\]
the ideal
\[
(s-b,t-a)
\]
or equivalently
\[
(t-a,s-b)
\]
is the ideal of the point
\[
(b,a)
\]
in \((s,t)\)-coordinates.

\subsubsection{Final conclusion}

So the correct interpretation is:

\begin{itemize}
	\item the indeterminates in \(k[s,t]\) are \(s\) and \(t\),
	\item the ideal \((t-a,s-b)\) imposes
	\[
	t=a,\qquad s=b,
	\]
	\item therefore it corresponds to the single point
	\[
	(s,t)=(b,a).
	\]
\end{itemize}

In particular,
\[
(t-a,s-b)
\]
does not describe two points. It describes one point in the affine plane with coordinates \((s,t)\).

\subsection{The ideal \((s-a,t-b)\) versus the quotient \(k[s,t]/(s-a,t-b)\)}

It is important to distinguish carefully between the ideal
\[
(s-a,t-b)\subseteq k[s,t]
\]
and the quotient ring
\[
k[s,t]/(s-a,t-b).
\]

These are not the same object.

\subsubsection{What the ideal \((s-a,t-b)\) consists of}

By definition, the ideal generated by \(s-a\) and \(t-b\) is
\[
(s-a,t-b)
=
\{(s-a)f(s,t)+(t-b)g(s,t)\;:\; f(s,t),g(s,t)\in k[s,t]\}.
\]

So an element of the ideal is any polynomial of the form
\[
(s-a)f(s,t)+(t-b)g(s,t).
\]

Thus the ideal consists of all finite \(k[s,t]\)-linear combinations of the generators \(s-a\) and \(t-b\).

In particular, it is not only made of terms of the form \((s-a)f\) alone or \((t-b)g\) alone, but of all sums
\[
(s-a)f+(t-b)g.
\]

\subsubsection{What the quotient \(k[s,t]/(s-a,t-b)\) consists of}

The quotient ring
\[
k[s,t]/(s-a,t-b)
\]
does \emph{not} consist of those sums themselves.  
Instead, it consists of equivalence classes of polynomials modulo the ideal \((s-a,t-b)\).

That is, two polynomials \(p(s,t)\) and \(q(s,t)\) define the same element in the quotient if and only if
\[
p(s,t)-q(s,t)\in (s-a,t-b).
\]

Equivalently,
\[
p \equiv q \pmod{(s-a,t-b)}.
\]

So the quotient identifies two polynomials whenever their difference lies in the ideal.

\subsubsection{Why the quotient is isomorphic to \(k\)}

In the quotient ring
\[
k[s,t]/(s-a,t-b),
\]
we impose the relations
\[
s=a,\qquad t=b.
\]
Therefore every polynomial \(f(s,t)\) becomes equal to the scalar
\[
f(a,b)\in k.
\]

Equivalently, consider the evaluation map
\[
\varphi:k[s,t]\longrightarrow k,
\qquad
f(s,t)\longmapsto f(a,b).
\]
This is a surjective ring homomorphism, and its kernel is exactly
\[
\ker(\varphi)=(s-a,t-b).
\]
Hence, by the First Isomorphism Theorem,
\[
k[s,t]/(s-a,t-b)\cong k.
\]

So the quotient ring is the field \(k\), not the ideal itself.

\subsubsection{A concrete example}

Take \(k=\mathbb{R}\), \(a=1\), and \(b=2\). Then
\[
(s-1,t-2)\subseteq \mathbb{R}[s,t]
\]
is the ideal consisting of all polynomials of the form
\[
(s-1)f(s,t)+(t-2)g(s,t).
\]

For example, the following elements belong to the ideal:
\[
(s-1)(s+t),\qquad (t-2)(3s^2),\qquad (s-1)(s+t)+(t-2)(3s^2).
\]

Now consider the quotient
\[
\mathbb{R}[s,t]/(s-1,t-2).
\]
In this quotient we have
\[
s=1,\qquad t=2.
\]
So for instance:
\[
[s^2+t]=[1^2+2]=[3],
\]
\[
[st+5]=[1\cdot 2+5]=[7].
\]

Thus every class in the quotient is represented by a real number, and the quotient is isomorphic to \(\mathbb{R}\).

\subsubsection{A useful characterization of the ideal}

A polynomial \(h(s,t)\in k[s,t]\) lies in the ideal \((s-a,t-b)\) if and only if
\[
h(a,b)=0.
\]

So the ideal \((s-a,t-b)\) is precisely the set of all polynomials vanishing at the point \((a,b)\).

This is another way to understand why the quotient by this ideal corresponds to evaluating polynomials at the point \((a,b)\).

\subsubsection{Final summary}

Thus:

\begin{itemize}
	\item the ideal
	\[
	(s-a,t-b)
	\]
	consists of all polynomials of the form
	\[
	(s-a)f(s,t)+(t-b)g(s,t),
	\]
	
	\item the quotient
	\[
	k[s,t]/(s-a,t-b)
	\]
	consists of equivalence classes of polynomials modulo that ideal,
	
	\item and this quotient is isomorphic to \(k\), because quotienting by \((s-a,t-b)\) means evaluating at the single point \((a,b)\).
\end{itemize}

\subsection{What it means that \(k[s,t]/(s-a,t-b)\cong k\)}

When we write
\[
k[s,t]/(s-a,t-b)\cong k,
\]
the idea is that in the quotient, every polynomial is identified only by its value at the point
\[
(a,b).
\]

So, informally, one may say that the quotient is ``the values of polynomials at the point \((a,b)\).''

\subsubsection{The precise meaning}

More precisely, the quotient is not merely a set of values in some loose sense.  
It is a ring, and indeed a \(k\)-algebra, which is isomorphic to the field \(k\).

The isomorphism is given by the evaluation map
\[
\Phi: k[s,t]/(s-a,t-b)\longrightarrow k,
\qquad
[f(s,t)]\longmapsto f(a,b).
\]

This is well defined because if
\[
[f(s,t)]=[g(s,t)]
\]
in the quotient, then
\[
f(s,t)-g(s,t)\in (s-a,t-b),
\]
and therefore
\[
f(a,b)=g(a,b).
\]

Thus the class of a polynomial depends only on its value at the point \((a,b)\).

\subsubsection{Why this is an isomorphism}

The map
\[
[f(s,t)]\mapsto f(a,b)
\]
preserves addition and multiplication:
\[
[f]+[g]=[f+g]
\quad\mapsto\quad
f(a,b)+g(a,b),
\]
and
\[
[f]\cdot[g]=[fg]
\quad\mapsto\quad
f(a,b)g(a,b).
\]

So the quotient has exactly the same ring structure as \(k\).  
That is why one says that
\[
k[s,t]/(s-a,t-b)\cong k.
\]

\subsubsection{Interpretation}

Thus the correct interpretation is:

\begin{itemize}
	\item the quotient remembers only the value of a polynomial at the point \((a,b)\),
	\item each residue class corresponds to one scalar in \(k\),
	\item and the ring operations in the quotient correspond exactly to the usual addition and multiplication in \(k\).
\end{itemize}

So when one says that
\[
k[s,t]/(s-a,t-b)\cong k,
\]
one means that the quotient is isomorphic to the field of values of polynomials at the point \((a,b)\), and this field is exactly \(k\).

\subsubsection{Example}

Take \(k=\mathbb{R}\), \(a=1\), \(b=2\). Then
\[
\mathbb{R}[s,t]/(s-1,t-2)\cong \mathbb{R},
\]
and the class of a polynomial is determined by its value at \((1,2)\).

For example:
\[
[s^2+t]\mapsto 1^2+2=3,
\]
\[
[s+t]\mapsto 1+2=3,
\]
so
\[
[s^2+t]=[s+t]
\]
in the quotient, because both polynomials have the same value at \((1,2)\).

\subsubsection{Final summary}

Therefore, saying that
\[
k[s,t]/(s-a,t-b)\cong k
\]
means that after imposing
\[
s=a,\qquad t=b,
\]
every polynomial reduces to its scalar value at that point, and the quotient ring is exactly the field \(k\) of those values.

\subsection{What the quotient by \(\bigl((s-a)(t-b)\bigr)\) means}

Now consider the principal ideal
\[
\bigl((s-a)(t-b)\bigr)\subseteq k[s,t].
\]
This is very different from the ideal
\[
(s-a,t-b).
\]

The ideal \((s-a,t-b)\) forces
\[
s=a,\qquad t=b,
\]
so it corresponds to a single point.

By contrast, the ideal
\[
\bigl((s-a)(t-b)\bigr)
\]
forces only the relation
\[
(s-a)(t-b)=0.
\]
It does \emph{not} force \(s=a\) and it does \emph{not} force \(t=b\) separately.

\subsubsection{What the ideal \(\bigl((s-a)(t-b)\bigr)\) consists of}

Since this is a principal ideal, it consists of all multiples of the polynomial \((s-a)(t-b)\):
\[
\bigl((s-a)(t-b)\bigr)
=
\left\{(s-a)(t-b)\,h(s,t)\;:\; h(s,t)\in k[s,t]\right\}.
\]

So an element of this ideal is any polynomial divisible by \((s-a)(t-b)\).

\subsubsection{What the quotient means}

The quotient ring
\[
k[s,t]\big/\bigl((s-a)(t-b)\bigr)
\]
consists of residue classes of polynomials modulo the relation
\[
(s-a)(t-b)=0.
\]

Thus in the quotient we have
\[
[s-a]\cdot [t-b]=0,
\]
but in general neither \([s-a]\) nor \([t-b]\) is zero.

So this quotient is not a field. In fact, it has zero divisors.

\subsubsection{A first important contrast with the point ideal}

For the point ideal \((s-a,t-b)\), one has
\[
k[s,t]/(s-a,t-b)\cong k,
\]
because every polynomial is reduced to its value at the single point \((a,b)\).

For the product ideal \(\bigl((s-a)(t-b)\bigr)\), this is no longer true.  
We do \emph{not} reduce to one scalar value.

Instead, the quotient remembers functions on the union of the two lines
\[
s=a
\qquad\text{and}\qquad
t=b.
\]

\subsubsection{Geometric meaning}

The zero set of the polynomial \((s-a)(t-b)\) is
\[
V\bigl((s-a)(t-b)\bigr)=V(s-a)\cup V(t-b).
\]
So geometrically this is the union of:

\begin{itemize}
	\item the vertical line
	\[
	s=a,
	\]
	\item the horizontal line
	\[
	t=b.
	\]
\end{itemize}

These two lines meet at the point
\[
(a,b).
\]

Thus
\[
k[s,t]\big/\bigl((s-a)(t-b)\bigr)
\]
is the coordinate ring of two intersecting lines, not of a single point.

\subsubsection{How to interpret elements of the quotient}

A polynomial \(f(s,t)\) determines:

\begin{itemize}
	\item its restriction to the line \(s=a\), namely
	\[
	f(a,t)\in k[t],
	\]
	\item its restriction to the line \(t=b\), namely
	\[
	f(s,b)\in k[s].
	\]
\end{itemize}

So instead of getting one scalar, we get two pieces of data:
\[
f(a,t)\quad\text{and}\quad f(s,b).
\]

These are not arbitrary, because they must agree at the intersection point \((a,b)\). Indeed,
\[
f(a,b)=f(a,b).
\]
So the two restrictions satisfy the compatibility condition
\[
f(a,t)\big|_{t=b}=f(s,b)\big|_{s=a}.
\]

This leads to the correct description of the quotient.

\subsubsection{The correct isomorphism}

Define a map
\[
\Phi:
k[s,t]\big/\bigl((s-a)(t-b)\bigr)
\longrightarrow
\left\{(u(t),v(s))\in k[t]\times k[s]\;:\; u(b)=v(a)\right\}
\]
by
\[
[f(s,t)]\longmapsto \bigl(f(a,t),\,f(s,b)\bigr).
\]

Then
\[
\Phi
\]
is a ring isomorphism.

So we have
\[
k[s,t]\big/\bigl((s-a)(t-b)\bigr)
\cong
\left\{(u(t),v(s))\in k[t]\times k[s]\;:\; u(b)=v(a)\right\}.
\]

This means that an element of the quotient is exactly a pair of polynomial functions:

\begin{itemize}
	\item one on the line \(s=a\),
	\item one on the line \(t=b\),
\end{itemize}

which agree at the crossing point \((a,b)\).

\subsubsection{Why this is the right description}

Let us explain why the previous map is correct.

\paragraph{Well definedness}

If two polynomials differ by a multiple of \((s-a)(t-b)\), then their restrictions to both lines are the same, because
\[
(s-a)(t-b)\big|_{s=a}=0,
\qquad
(s-a)(t-b)\big|_{t=b}=0.
\]
So the pair \(\bigl(f(a,t),f(s,b)\bigr)\) depends only on the residue class of \(f\).

\paragraph{Kernel}

Suppose
\[
\Phi([f])=(0,0).
\]
Then
\[
f(a,t)=0
\quad\text{and}\quad
f(s,b)=0.
\]

From \(f(a,t)=0\), it follows that \(f\) is divisible by \(s-a\), so
\[
f=(s-a)g
\]
for some \(g\in k[s,t]\).

Now substituting \(t=b\), we get
\[
0=f(s,b)=(s-a)g(s,b).
\]
Since \(k[s]\) is an integral domain and \(s-a\neq 0\), we conclude
\[
g(s,b)=0.
\]
Hence \(g\) is divisible by \(t-b\), so
\[
g=(t-b)h
\]
for some \(h\in k[s,t]\).

Therefore
\[
f=(s-a)(t-b)h,
\]
so
\[
f\in \bigl((s-a)(t-b)\bigr).
\]

Thus the kernel is exactly
\[
\ker(\Phi)=\bigl((s-a)(t-b)\bigr).
\]

\paragraph{Surjectivity}

Now take a pair
\[
(u(t),v(s))\in k[t]\times k[s]
\]
such that
\[
u(b)=v(a).
\]
Set
\[
c:=u(b)=v(a),
\]
and define
\[
f(s,t):=u(t)+v(s)-c.
\]
Then
\[
f(a,t)=u(t)+v(a)-c=u(t),
\]
and
\[
f(s,b)=u(b)+v(s)-c=v(s).
\]
So every compatible pair comes from some polynomial \(f\).

Hence \(\Phi\) is surjective.

\subsubsection{Why this quotient is not \(k\)}

The quotient
\[
k[s,t]\big/\bigl((s-a)(t-b)\bigr)
\]
is not isomorphic to \(k\), because it has nonzero zero divisors. Indeed, in the quotient,
\[
[s-a]\neq 0,\qquad [t-b]\neq 0,
\]
but
\[
[s-a]\cdot [t-b]=0.
\]

A field cannot have nonzero zero divisors, so this quotient is not a field.

Thus it cannot be isomorphic to \(k\).

\subsubsection{Why this quotient is not \(k[t]\times k[s]\)}

It is also important to note that
\[
k[s,t]\big/\bigl((s-a)(t-b)\bigr)
\]
is not the full direct product
\[
k[t]\times k[s].
\]

The reason is that the two lines
\[
s=a
\quad\text{and}\quad
t=b
\]
are not disjoint: they intersect at \((a,b)\).

If we had the full product \(k[t]\times k[s]\), then the functions on the two lines would be completely independent. But in reality they must agree at the intersection point.

That is why the correct description is the subring
\[
\left\{(u(t),v(s))\in k[t]\times k[s]\;:\; u(b)=v(a)\right\},
\]
not the full product.

\subsubsection{Concrete example}

Take \(k=\mathbb{R}\), \(a=1\), \(b=2\). Then
\[
\mathbb{R}[s,t]\big/\bigl((s-1)(t-2)\bigr)
\]
is the coordinate ring of the union of the two lines
\[
s=1
\qquad\text{and}\qquad
t=2.
\]

In this quotient, we have
\[
[s-1]\cdot [t-2]=0,
\]
but neither factor is zero.

A polynomial \(f(s,t)\) determines:

\begin{itemize}
	\item the polynomial \(f(1,t)\) on the line \(s=1\),
	\item the polynomial \(f(s,2)\) on the line \(t=2\),
\end{itemize}
and these two must agree at the point \((1,2)\).

For example, let
\[
f(s,t)=s+t.
\]
Then
\[
f(1,t)=1+t,
\qquad
f(s,2)=s+2.
\]
These agree at the intersection:
\[
f(1,2)=3.
\]

So the class of \(s+t\) in the quotient corresponds to the compatible pair
\[
(1+t,\; s+2).
\]

\subsubsection{Comparison with the single-point case}

It is useful to compare the two quotients:

\[
k[s,t]/(s-a,t-b)\cong k,
\]
because this is the coordinate ring of one point.

But
\[
k[s,t]\big/\bigl((s-a)(t-b)\bigr)
\]
is the coordinate ring of the union of two intersecting lines, so it is much larger and has zero divisors.

Thus:

\begin{itemize}
	\item \((s-a,t-b)\) corresponds to one point,
	\item \(\bigl((s-a)(t-b)\bigr)\) corresponds to two lines crossing at that point.
\end{itemize}

\subsubsection{Final summary}

The quotient
\[
k[s,t]\big/\bigl((s-a)(t-b)\bigr)
\]
does not mean ``evaluate at one point.'' Instead, it means:

\begin{itemize}
	\item keep track of polynomial functions on the line \(s=a\),
	\item keep track of polynomial functions on the line \(t=b\),
	\item require that these two functions agree at the intersection point \((a,b)\).
\end{itemize}

Equivalently,
\[
k[s,t]\big/\bigl((s-a)(t-b)\bigr)
\cong
\left\{(u(t),v(s))\in k[t]\times k[s]\;:\; u(b)=v(a)\right\}.
\]

So the quotient by \(\bigl((s-a)(t-b)\bigr)\) describes the union of two intersecting lines, not a single value in \(k\).


\subsection{Distinct principal ideals and their quotient rings for unions of two lines}

Consider the two polynomials
\[
f_1=(s+t-1)(s-2t+3),
\qquad
f_2=(s-5)(s+t)
\]
in the polynomial ring
\[
k[s,t],
\]
where \(k\) is a field.

A natural question is:

\medskip

Do these define distinct ideals?
Do they give distinct quotient rings?
And in what sense are the corresponding geometric objects the same or different?

\subsubsection{The ideals are distinct}

Yes, the ideals
\[
(f_1)
\qquad\text{and}\qquad
(f_2)
\]
are distinct ideals of \(k[s,t]\).

The reason is the following general fact.

\begin{proposition}
	Let \(R\) be an integral domain, and let \(f,g\in R\) be nonzero. Then
	\[
	(f)=(g)
	\]
	if and only if there exists a unit \(u\in R^\times\) such that
	\[
	f=ug.
	\]
	In other words, two nonzero principal ideals are equal if and only if their generators are associates.
\end{proposition}

\begin{proof}
	Assume first that
	\[
	(f)=(g).
	\]
	Then \(f\in (g)\), so there exists \(a\in R\) such that
	\[
	f=ag.
	\]
	Similarly, \(g\in (f)\), so there exists \(b\in R\) such that
	\[
	g=bf.
	\]
	Substituting the second equation into the first gives
	\[
	f=a(bf)=(ab)f.
	\]
	Since \(R\) is an integral domain and \(f\neq 0\), we may cancel \(f\) and obtain
	\[
	ab=1.
	\]
	Hence \(a\) is a unit, and
	\[
	f=ag
	\]
	shows that \(f\) and \(g\) differ by multiplication by a unit.
	
	Conversely, if
	\[
	f=ug
	\]
	for a unit \(u\), then clearly
	\[
	(f)=(g),
	\]
	since multiplication by a unit does not change the principal ideal.
\end{proof}

Now apply this to the ring
\[
k[s,t].
\]
Its units are exactly the nonzero scalars in \(k\). Therefore
\[
(f_1)=(f_2)
\]
would force
\[
f_1=c\,f_2
\qquad\text{for some } c\in k^\times.
\]
But this is impossible, because the factorizations are different:

\[
f_1=(s+t-1)(s-2t+3),
\qquad
f_2=(s-5)(s+t),
\]
and \(f_1\) is not a nonzero scalar multiple of \(f_2\).

Hence
\[
(f_1)\neq (f_2).
\]

\subsubsection{Geometric meaning of the two polynomials}

Each polynomial is a product of two linear factors, so each defines a union of two lines in the affine plane \(\mathbb{A}^2_k\).

For the first polynomial,
\[
f_1=(s+t-1)(s-2t+3),
\]
the zero locus is
\[
V(f_1)=V(s+t-1)\cup V(s-2t+3).
\]
So this is the union of the two lines
\[
s+t-1=0
\qquad\text{and}\qquad
s-2t+3=0.
\]

For the second polynomial,
\[
f_2=(s-5)(s+t),
\]
the zero locus is
\[
V(f_2)=V(s-5)\cup V(s+t).
\]
So this is the union of the two lines
\[
s-5=0
\qquad\text{and}\qquad
s+t=0.
\]

Thus the two polynomials define two different pairs of lines, hence two different subsets of the affine plane.

\subsubsection{Each pair of lines intersects in one point}

In the first case, the two lines are
\[
s+t-1=0,
\qquad
s-2t+3=0.
\]
Subtracting the equations gives
\[
3t-4=0,
\qquad
t=\frac{4}{3},
\]
and then
\[
s=1-t=1-\frac{4}{3}=-\frac{1}{3}.
\]
So the two lines intersect at
\[
\left(-\frac13,\frac43\right).
\]

In the second case, the two lines are
\[
s-5=0,
\qquad
s+t=0.
\]
From \(s=5\), the second equation gives
\[
5+t=0,
\qquad
t=-5.
\]
So the intersection point is
\[
(5,-5).
\]

Hence in both examples we have a union of two distinct lines meeting transversely in one point.

\subsubsection{The quotient rings come from different ideals}

Since the ideals are different, the quotients
\[
k[s,t]/(f_1)
\qquad\text{and}\qquad
k[s,t]/(f_2)
\]
are quotients by different ideals inside the same polynomial ring.

So as \emph{presentations inside \(k[s,t]\)}, they are different quotient constructions.

Equivalently, the coordinate rings correspond to different closed subschemes of \(\mathbb{A}^2_k\).

\subsubsection{But the quotient rings are abstractly isomorphic}

Although the ideals are distinct inside \(k[s,t]\), the quotient rings are abstractly isomorphic as \(k\)-algebras.

The reason is that each quotient is the coordinate ring of a union of two distinct intersecting lines, and any such configuration can be transformed by an affine change of coordinates into the standard crossing
\[
uv=0.
\]

Thus each quotient is isomorphic to
\[
k[u,v]/(uv).
\]

\subsubsection{General principle}

Suppose
\[
\ell_1(s,t),\ell_2(s,t)\in k[s,t]
\]
are two distinct affine linear polynomials whose linear parts are linearly independent. Then the two lines
\[
\ell_1=0,
\qquad
\ell_2=0
\]
intersect in one point.

In that case, the map
\[
\Phi:k[u,v]\longrightarrow k[s,t]
\]
defined by
\[
u\longmapsto \ell_1(s,t),
\qquad
v\longmapsto \ell_2(s,t)
\]
is an isomorphism of \(k\)-algebras.

Indeed, since the linear parts of \(\ell_1\) and \(\ell_2\) are linearly independent, the corresponding \(2\times 2\) coefficient matrix is invertible, so one can solve for \(s\) and \(t\) as affine linear polynomials in \(u\) and \(v\). Hence \(\Phi\) has an inverse.

Therefore
\[
k[s,t]/(\ell_1\ell_2)\cong k[u,v]/(uv).
\]

This shows that the coordinate ring of any pair of distinct intersecting lines is abstractly isomorphic to the standard nodal ring
\[
k[u,v]/(uv).
\]

\subsubsection{Applying this to the first example}

Set
\[
u=s+t-1,
\qquad
v=s-2t+3.
\]
Then
\[
f_1=uv.
\]

The linear coefficient matrix is
\[
\begin{pmatrix}
	1 & 1\\
	1 & -2
\end{pmatrix},
\]
whose determinant is
\[
-2-1=-3.
\]
Assuming the characteristic of \(k\) does not make this determinant vanish, the matrix is invertible, and \(u,v\) form a new affine coordinate system.

So
\[
k[s,t]/\bigl((s+t-1)(s-2t+3)\bigr)\cong k[u,v]/(uv).
\]

If one wants an explicit inverse, solve:
\[
u=s+t-1,
\qquad
v=s-2t+3.
\]
Subtracting gives
\[
u-v=3t-4,
\]
so
\[
t=\frac{u-v+4}{3}.
\]
Then
\[
s=u-t+1
= u-\frac{u-v+4}{3}+1
= \frac{2u+v-1}{3}.
\]

Thus \(s\) and \(t\) are affine linear expressions in \(u\) and \(v\), confirming the isomorphism.

\subsubsection{Applying this to the second example}

Set
\[
u=s-5,
\qquad
v=s+t.
\]
Then
\[
f_2=uv.
\]

The linear coefficient matrix is
\[
\begin{pmatrix}
	1 & 0\\
	1 & 1
\end{pmatrix},
\]
whose determinant is
\[
1.
\]
So this change of variables is invertible over every field.

Hence
\[
k[s,t]/\bigl((s-5)(s+t)\bigr)\cong k[u,v]/(uv).
\]

Indeed, from
\[
u=s-5,
\qquad
v=s+t,
\]
we get
\[
s=u+5,
\qquad
t=v-s=v-u-5.
\]

So again \(s\) and \(t\) are affine linear expressions in \(u\) and \(v\), and the quotient is isomorphic to the standard ring
\[
k[u,v]/(uv).
\]

\subsubsection{Conclusion about the quotients}

Therefore, the correct conclusion is:

\begin{itemize}
	\item the ideals
	\[
	\bigl((s+t-1)(s-2t+3)\bigr)
	\quad\text{and}\quad
	\bigl((s-5)(s+t)\bigr)
	\]
	are distinct ideals of \(k[s,t]\),
	
	\item the quotient rings
	\[
	k[s,t]/\bigl((s+t-1)(s-2t+3)\bigr)
	\quad\text{and}\quad
	k[s,t]/\bigl((s-5)(s+t)\bigr)
	\]
	are quotients by different ideals,
	
	\item but abstractly, both quotient rings are isomorphic to
	\[
	k[u,v]/(uv),
	\]
	so they are isomorphic to each other as \(k\)-algebras.
\end{itemize}

\subsubsection{Why this does not contradict distinctness of the ideals}

There is no contradiction in saying that two ideals are distinct but the quotient rings are isomorphic.

Indeed, these are two different statements:

\begin{itemize}
	\item saying
	\[
	I\neq J
	\]
	means the ideals are different subsets of the ring \(k[s,t]\),
	
	\item saying
	\[
	k[s,t]/I \cong k[s,t]/J
	\]
	means the quotient rings are abstractly isomorphic as rings.
\end{itemize}

This is analogous to geometry: two different curves drawn in different positions in the plane may still be abstractly the same type of curve.

Here, the two unions of lines sit in different places and with different equations inside \(\mathbb{A}^2_k\), but each has the same local and global algebraic structure: two distinct lines crossing in one point.

\subsubsection{Relation with the earlier discussion}

This fits exactly with the earlier discussion of quotients of the form
\[
k[s,t]/\bigl((s-a)(t-b)\bigr).
\]
Such a quotient is the coordinate ring of the union of the two coordinate-type lines
\[
s=a
\qquad\text{and}\qquad
t=b.
\]
It is not a field, because it has zero divisors:
\[
[s-a]\neq 0,
\qquad
[t-b]\neq 0,
\qquad
[s-a][t-b]=0.
\]

Likewise, the quotients
\[
k[s,t]/\bigl((s+t-1)(s-2t+3)\bigr)
\qquad\text{and}\qquad
k[s,t]/\bigl((s-5)(s+t)\bigr)
\]
are also coordinate rings of unions of two intersecting lines, so they again have zero divisors and are again isomorphic to the standard crossing ring
\[
k[u,v]/(uv).
\]

So the earlier intuition remains correct:

\begin{itemize}
	\item quotient by \((s-a,t-b)\): one point, hence quotient is \(k\),
	\item quotient by \(\bigl((s-a)(t-b)\bigr)\): two intersecting lines, hence quotient is a nodal ring with zero divisors,
	\item quotient by any product of two distinct intersecting linear factors: same abstract ring type as \(k[u,v]/(uv)\).
\end{itemize}

\subsubsection{Final summary}

Let
\[
f_1=(s+t-1)(s-2t+3),
\qquad
f_2=(s-5)(s+t).
\]
Then:

\begin{enumerate}
	\item
	\[
	(f_1)\neq (f_2)
	\]
	as ideals in \(k[s,t]\),
	
	\item the quotient constructions
	\[
	k[s,t]/(f_1)
	\qquad\text{and}\qquad
	k[s,t]/(f_2)
	\]
	come from different ideals and thus define different embedded subschemes of \(\mathbb{A}^2_k\),
	
	\item nevertheless,
	\[
	k[s,t]/(f_1)\cong k[u,v]/(uv)\cong k[s,t]/(f_2),
	\]
	so abstractly they are isomorphic as \(k\)-algebras.
\end{enumerate}

Thus the ideals are different, the quotient presentations are different, but the quotient rings have the same abstract algebraic structure.

\subsection{Irreducible quadratic factors over \(\mathbb{R}\) and why the quotient is not the same as for real linear factors}

A natural question is whether, over the field \(\mathbb{R}\), one may replace linear factors by irreducible quadratic factors with no real roots and still obtain ``the same kind'' of quotient ring as in the case of products of real linear factors.

The answer is, in general, \textbf{no}.

\subsubsection{The basic comparison in one variable}

First compare the two quotient rings
\[
\mathbb{R}[t]/\bigl((t-1)(t-2)\bigr)
\qquad\text{and}\qquad
\mathbb{R}[t]/\bigl((t^2+1)(t^2+4)\bigr).
\]

These are very different as \(\mathbb{R}\)-algebras.

\paragraph{Case 1: two real linear factors.}

Since the ideals \((t-1)\) and \((t-2)\) are comaximal, the Chinese Remainder Theorem gives
\[
\mathbb{R}[t]/\bigl((t-1)(t-2)\bigr)
\cong
\mathbb{R}[t]/(t-1)\times \mathbb{R}[t]/(t-2).
\]
But
\[
\mathbb{R}[t]/(t-1)\cong \mathbb{R},
\qquad
\mathbb{R}[t]/(t-2)\cong \mathbb{R},
\]
so
\[
\mathbb{R}[t]/\bigl((t-1)(t-2)\bigr)\cong \mathbb{R}\times \mathbb{R}.
\]

Thus the quotient corresponds to two distinct real points.

\paragraph{Case 2: two irreducible quadratic factors.}

Now consider
\[
\mathbb{R}[t]/\bigl((t^2+1)(t^2+4)\bigr).
\]
Again the two factors are comaximal, because
\[
(t^2+4)-(t^2+1)=3,
\]
and \(3\) is a unit in \(\mathbb{R}[t]\). Hence by the Chinese Remainder Theorem,
\[
\mathbb{R}[t]/\bigl((t^2+1)(t^2+4)\bigr)
\cong
\mathbb{R}[t]/(t^2+1)\times \mathbb{R}[t]/(t^2+4).
\]

Now both quotients are fields isomorphic to \(\mathbb{C}\):
\[
\mathbb{R}[t]/(t^2+1)\cong \mathbb{C},
\qquad
\mathbb{R}[t]/(t^2+4)\cong \mathbb{C}.
\]
Indeed, in the second case one may send the class of \(t\) to \(2i\).

Therefore
\[
\mathbb{R}[t]/\bigl((t^2+1)(t^2+4)\bigr)\cong \mathbb{C}\times \mathbb{C}.
\]

So this quotient corresponds to two closed points whose residue fields are \(\mathbb{C}\), not \(\mathbb{R}\).

\subsubsection{Why \(\mathbb{R}\times\mathbb{R}\) and \(\mathbb{C}\times\mathbb{C}\) are not the same}

These rings are not isomorphic as \(\mathbb{R}\)-algebras.

One easy reason is that \(\mathbb{C}\times \mathbb{C}\) contains an element whose square is \((-1,-1)\), namely
\[
(i,i)^2=(-1,-1),
\]
so it contains an element whose square is the image of \(-1\).

By contrast, in \(\mathbb{R}\times \mathbb{R}\), if
\[
(a,b)^2=(-1,-1),
\]
then we would need
\[
a^2=-1,\qquad b^2=-1,
\]
which is impossible in \(\mathbb{R}\).

Hence
\[
\mathbb{R}\times \mathbb{R}\not\cong \mathbb{C}\times \mathbb{C}
\]
as \(\mathbb{R}\)-algebras.

So the quotient coming from two irreducible quadratic factors is not the same as the quotient coming from two distinct real linear factors.

\subsubsection{Scheme-theoretic meaning}

This is an important point in scheme theory.

Even though the polynomial
\[
t^2+1
\]
has no real zero in the usual sense, the scheme
\[
\Spec\bigl(\mathbb{R}[t]/(t^2+1)\bigr)
\]
is not empty. It consists of one closed point whose residue field is
\[
\mathbb{R}[t]/(t^2+1)\cong \mathbb{C}.
\]

Likewise,
\[
\Spec\bigl(\mathbb{R}[t]/(t^2+4)\bigr)
\]
is another closed point with residue field \(\mathbb{C}\).

Thus
\[
\Spec\!\left(\mathbb{R}[t]/\bigl((t^2+1)(t^2+4)\bigr)\right)
\]
is the disjoint union of two closed points, each having residue field \(\mathbb{C}\).

By contrast,
\[
\Spec\!\left(\mathbb{R}[t]/\bigl((t-1)(t-2)\bigr)\right)
\]
is the disjoint union of two real points, each having residue field \(\mathbb{R}\).

So even though both quotient rings split as products of two fields, the fields are different.

\subsubsection{What happens in two variables}

Now consider a polynomial in two variables such as
\[
(s^2+t^2+1)(s^2+t^2+4).
\]
Neither factor has any real zero, since
\[
s^2+t^2+1>0
\qquad\text{and}\qquad
s^2+t^2+4>0
\]
for all \(s,t\in \mathbb{R}\).

Nevertheless, the quotient ring
\[
\mathbb{R}[s,t]\Big/\bigl((s^2+t^2+1)(s^2+t^2+4)\bigr)
\]
is perfectly meaningful algebraically.

As before, the two factors are comaximal, since
\[
(s^2+t^2+4)-(s^2+t^2+1)=3.
\]
Therefore
\[
\mathbb{R}[s,t]\Big/\bigl((s^2+t^2+1)(s^2+t^2+4)\bigr)
\cong
\mathbb{R}[s,t]/(s^2+t^2+1)\times \mathbb{R}[s,t]/(s^2+t^2+4).
\]

So this quotient ring splits into a product of two coordinate rings of irreducible real conics with no real points.

This is very different from the ring
\[
\mathbb{R}[u,v]/(uv),
\]
which is the coordinate ring of two real lines crossing at the origin.

\subsubsection{Why this is different from the crossing-lines ring}

The ring
\[
\mathbb{R}[u,v]/(uv)
\]
has zero divisors:
\[
[u]\neq 0,\qquad [v]\neq 0,\qquad [u][v]=0.
\]
Geometrically, it describes a reducible real curve consisting of two real lines intersecting in one real point.

By contrast,
\[
\mathbb{R}[s,t]\Big/\bigl((s^2+t^2+1)(s^2+t^2+4)\bigr)
\]
is also reducible, because it is a product of two nonconstant factors, but its components are not real lines and in fact have no real points at all.

So the geometry is entirely different:

\begin{itemize}
	\item
	\[
	uv=0
	\]
	describes two real coordinate axes crossing in a real point;
	
	\item
	\[
	(s^2+t^2+1)(s^2+t^2+4)=0
	\]
	describes two irreducible real conics, neither of which has a real point.
\end{itemize}

Thus the corresponding quotient rings are not of the same geometric type over \(\mathbb{R}\).

\subsubsection{A useful general principle}

Over \(\mathbb{R}\), the nature of the irreducible factors matters.

\begin{itemize}
	\item Distinct real linear factors give components with residue field \(\mathbb{R}\).
	\item Irreducible quadratic factors give closed points or components whose residue fields are extensions of \(\mathbb{R}\), typically \(\mathbb{C}\).
\end{itemize}

So even if the formal expression is a product of two irreducible factors, the resulting quotient ring depends strongly on what kind of factors they are.

In particular, products of irreducible quadratics over \(\mathbb{R}\) are generally \emph{not} algebraically the same as products of distinct real linear factors.

\subsubsection{A further comparison via residue fields}

The difference can also be detected through residue fields.

For
\[
\mathbb{R}[t]/\bigl((t-1)(t-2)\bigr)\cong \mathbb{R}\times \mathbb{R},
\]
the two factors have residue fields
\[
\mathbb{R},\qquad \mathbb{R}.
\]

For
\[
\mathbb{R}[t]/\bigl((t^2+1)(t^2+4)\bigr)\cong \mathbb{C}\times \mathbb{C},
\]
the two factors have residue fields
\[
\mathbb{C},\qquad \mathbb{C}.
\]

Since residue fields are intrinsic local invariants, these quotient rings cannot be isomorphic as \(\mathbb{R}\)-algebras.

\subsubsection{Final conclusion}

Therefore, over \(\mathbb{R}\), if one takes two irreducible quadratic factors with no real zeros, the resulting quotient ring is generally \emph{not} the same as the quotient coming from two distinct real linear factors.

For example,
\[
\mathbb{R}[t]/\bigl((t-1)(t-2)\bigr)\cong \mathbb{R}\times \mathbb{R},
\]
while
\[
\mathbb{R}[t]/\bigl((t^2+1)(t^2+4)\bigr)\cong \mathbb{C}\times \mathbb{C}.
\]

Likewise, in two variables,
\[
\mathbb{R}[u,v]/(uv)
\]
is the coordinate ring of two real lines crossing, whereas
\[
\mathbb{R}[s,t]\Big/\bigl((s^2+t^2+1)(s^2+t^2+4)\bigr)
\]
is the coordinate ring of two irreducible real components with no real points.

\medskip

\textbf{Solution.}
No, in general the quotient is not the same. Distinct real linear factors over \(\mathbb{R}\) lead to quotients such as
\[
\mathbb{R}\times \mathbb{R},
\]
whereas distinct irreducible quadratic factors with no real roots lead to quotients such as
\[
\mathbb{C}\times \mathbb{C}.
\]
These are not isomorphic as \(\mathbb{R}\)-algebras. In higher dimensions the difference persists: products of real linear factors describe unions of real hyperplanes, while products of irreducible quadratic factors may describe components with no real points and different residue fields.

\subsection{Why \(\mathbb{C}\) appears when we quotient \(\mathbb{R}[t]\) by \((t^2+1)\)}

A natural question is the following:

\medskip

Why do we suddenly obtain \(\mathbb{C}\) when starting from the real polynomial ring \(\mathbb{R}[t]\)?  
Why do we seem to ``shift'' from \(\mathbb{R}\) to \(\mathbb{C}\)?

\medskip

The answer is that we do \emph{not} arbitrarily change the base field. Rather, the quotient ring itself naturally becomes a field isomorphic to \(\mathbb{C}\).

\subsubsection{The quotient imposes the relation \(t^2=-1\)}

Consider the quotient ring
\[
\mathbb{R}[t]/(t^2+1).
\]
In this quotient, the polynomial \(t^2+1\) becomes zero, so the class of \(t\), denoted \([t]\), satisfies
\[
[t]^2+1=0.
\]
Equivalently,
\[
[t]^2=-1.
\]

Thus the quotient ring contains an element whose square is \(-1\).

\subsubsection{Why this cannot be \(\mathbb{R}\)}

In the field \(\mathbb{R}\), there is no element whose square is \(-1\).  
So the quotient
\[
\mathbb{R}[t]/(t^2+1)
\]
cannot possibly be isomorphic to \(\mathbb{R}\).

Instead, the quotient creates a new element, namely the class \([t]\), which behaves exactly like the imaginary unit \(i\).

This is why the quotient is naturally isomorphic to \(\mathbb{C}\).

\subsubsection{The precise isomorphism}

Define a ring homomorphism
\[
\varphi:\mathbb{R}[t]\longrightarrow \mathbb{C}
\]
by evaluation at \(i\):
\[
\varphi(f(t))=f(i).
\]

This map is surjective, because every complex number can be written in the form
\[
a+bi
\]
with \(a,b\in \mathbb{R}\), and this is the value of the polynomial
\[
a+bt
\]
at \(t=i\).

Now the kernel of \(\varphi\) consists of all polynomials \(f(t)\in \mathbb{R}[t]\) such that
\[
f(i)=0.
\]
Since \(t^2+1\) is the minimal polynomial of \(i\) over \(\mathbb{R}\), we get
\[
\ker(\varphi)=(t^2+1).
\]

Therefore, by the First Isomorphism Theorem,
\[
\mathbb{R}[t]/(t^2+1)\cong \mathbb{C}.
\]

So \(\mathbb{C}\) appears because the quotient ring itself is naturally isomorphic to \(\mathbb{C}\).

\subsubsection{What is really happening algebraically}

One may think of \(\mathbb{R}[t]\) as the ring of polynomials in a formal variable \(t\).  
When we quotient by \((t^2+1)\), we are imposing the rule
\[
t^2=-1.
\]

So the quotient is the smallest \(\mathbb{R}\)-algebra in which there exists an element whose square is \(-1\).

But that is exactly what \(\mathbb{C}\) is: the smallest field containing \(\mathbb{R}\) and an element \(i\) with
\[
i^2=-1.
\]

Hence
\[
\mathbb{R}[t]/(t^2+1)\cong \mathbb{C}.
\]

\subsubsection{Comparison with quotienting by a linear factor}

Compare this with the quotient
\[
\mathbb{R}[t]/(t-a),
\qquad a\in \mathbb{R}.
\]
In that quotient we impose
\[
t=a,
\]
which is already a relation realizable inside \(\mathbb{R}\).

Therefore
\[
\mathbb{R}[t]/(t-a)\cong \mathbb{R}.
\]

So:

\begin{itemize}
	\item quotienting by \((t-a)\) gives \(\mathbb{R}\), because \(a\in\mathbb{R}\) already exists in the ground field;
	\item quotienting by \((t^2+1)\) gives a larger field, because \(\mathbb{R}\) does not contain a root of \(t^2+1\).
\end{itemize}

\subsubsection{Scheme-theoretic interpretation}

From the scheme-theoretic point of view, the prime ideal
\[
(t^2+1)\subset \mathbb{R}[t]
\]
defines a closed point of
\[
\Spec(\mathbb{R}[t]).
\]

Its residue field is
\[
k(x)=\mathbb{R}[t]/(t^2+1)\cong \mathbb{C}.
\]

So this point is a closed point of the \(\mathbb{R}\)-scheme \(\Spec(\mathbb{R}[t])\), but it is not an \(\mathbb{R}\)-rational point, because its residue field is \(\mathbb{C}\), not \(\mathbb{R}\).

This is why in algebraic geometry one says that \((t^2+1)\) defines a point with complex residue field.

\subsubsection{Final summary}

Thus we do not arbitrarily ``switch'' from \(\mathbb{R}\) to \(\mathbb{C}\).  
Rather, the quotient ring
\[
\mathbb{R}[t]/(t^2+1)
\]
forces the existence of an element whose square is \(-1\), and the resulting field is naturally isomorphic to \(\mathbb{C}\).

In short:

\begin{itemize}
	\item the base ring is still \(\mathbb{R}[t]\),
	\item quotienting by \((t^2+1)\) imposes the relation
	\[
	t^2=-1,
	\]
	\item \(\mathbb{R}\) has no such element,
	\item so the quotient becomes a larger field,
	\item and that field is isomorphic to \(\mathbb{C}\).
\end{itemize}

\subsection{Polynomial rings, ideals, quotients, and tensor products}

A very useful next step is to understand how polynomial rings, ideals, quotients, and tensor products interact.  
This explains systematically why residue fields change, why complex numbers appear from real polynomial quotients, and what happens when one enlarges the base field.

\subsubsection{General philosophy}

Suppose \(A\) is a ring, \(I\subseteq A\) an ideal, and \(B\) an \(A\)-algebra.  
Then there are two basic constructions:

\begin{itemize}
	\item first form the quotient \(A/I\),
	\item then extend scalars from \(A\) to \(B\) by tensoring with \(B\).
\end{itemize}

The fundamental fact is that these two operations commute:
\[
(A/I)\otimes_A B \cong B/IB.
\]

This simple formula is one of the most important bridges between quotients and tensor products.

It says:

\begin{quote}
	If one first imposes the relations defining \(I\), and then changes the base ring from \(A\) to \(B\), one gets the same result as first changing scalars to \(B\), then imposing the transported relations.
\end{quote}

\subsubsection{Polynomial rings and change of scalars}

Let \(k\to K\) be a field extension. Then
\[
k[x_1,\dots,x_n]\otimes_k K \cong K[x_1,\dots,x_n].
\]

So tensoring a polynomial ring with a larger field just changes the coefficients.

For example:
\[
\mathbb{R}[t]\otimes_{\mathbb{R}} \mathbb{C}\cong \mathbb{C}[t].
\]

Likewise,
\[
k[s,t]\otimes_k K\cong K[s,t].
\]

Thus tensoring with a larger field does not change the variables; it only enlarges the coefficient field.

\subsubsection{Quotients and tensor products}

Now let \(f\in k[t]\). Then
\[
\bigl(k[t]/(f)\bigr)\otimes_k K
\cong
K[t]/(f),
\]
where the same polynomial \(f\) is now viewed as an element of \(K[t]\).

This is a special case of the general rule
\[
(A/I)\otimes_A B \cong B/IB.
\]

So if one begins with a quotient ring over \(k\), then tensors up to \(K\), one gets the quotient by the same equation over the larger field \(K\).

This is exactly what happens when one passes from \(\mathbb{R}\) to \(\mathbb{C}\).

\subsubsection{The fundamental example: \(\mathbb{R}[t]/(t^2+1)\)}

Consider
\[
A=\mathbb{R}[t]/(t^2+1).
\]
We already know that
\[
A\cong \mathbb{C}
\]
as an \(\mathbb{R}\)-algebra.

Now tensor with \(\mathbb{C}\) over \(\mathbb{R}\). Then
\[
A\otimes_{\mathbb{R}}\mathbb{C}
\cong
\bigl(\mathbb{R}[t]/(t^2+1)\bigr)\otimes_{\mathbb{R}}\mathbb{C}.
\]
Using the quotient-tensor rule,
\[
\bigl(\mathbb{R}[t]/(t^2+1)\bigr)\otimes_{\mathbb{R}}\mathbb{C}
\cong
\mathbb{C}[t]/(t^2+1).
\]

But over \(\mathbb{C}\), the polynomial factors:
\[
t^2+1=(t-i)(t+i).
\]
Hence
\[
\mathbb{C}[t]/(t^2+1)
\cong
\mathbb{C}[t]/\bigl((t-i)(t+i)\bigr).
\]

Since the ideals \((t-i)\) and \((t+i)\) are comaximal, the Chinese remainder theorem gives
\[
\mathbb{C}[t]/(t^2+1)\cong \mathbb{C}[t]/(t-i)\times \mathbb{C}[t]/(t+i)\cong \mathbb{C}\times \mathbb{C}.
\]

Therefore
\[
\mathbb{C}\otimes_{\mathbb{R}}\mathbb{C}\cong \mathbb{C}\times \mathbb{C}.
\]

This is a very important example.  
It shows that after extending scalars from \(\mathbb{R}\) to \(\mathbb{C}\), the single irreducible real factor \(t^2+1\) splits into two linear complex factors.

\subsubsection{Interpretation of the previous example}

Over \(\mathbb{R}\), the quotient
\[
\mathbb{R}[t]/(t^2+1)
\]
is one field, namely \(\mathbb{C}\), viewed as an \(\mathbb{R}\)-algebra.

Over \(\mathbb{C}\), the same equation becomes
\[
t^2+1=(t-i)(t+i),
\]
so the corresponding complexified quotient becomes
\[
\mathbb{C}\times \mathbb{C}.
\]

So tensoring with \(\mathbb{C}\) reveals that the one real closed point with residue field \(\mathbb{C}\) becomes two \(\mathbb{C}\)-rational points.

This is a basic example of how extension of scalars can split an irreducible object into simpler pieces.

\subsubsection{Another example: two irreducible quadratic factors over \(\mathbb{R}\)}

Consider
\[
A=\mathbb{R}[t]/\bigl((t^2+1)(t^2+4)\bigr).
\]
Over \(\mathbb{R}\), the Chinese remainder theorem gives
\[
A
\cong
\mathbb{R}[t]/(t^2+1)\times \mathbb{R}[t]/(t^2+4)
\cong
\mathbb{C}\times \mathbb{C}.
\]

Now tensor with \(\mathbb{C}\):
\[
A\otimes_{\mathbb{R}}\mathbb{C}
\cong
\mathbb{C}[t]/\bigl((t^2+1)(t^2+4)\bigr).
\]
Over \(\mathbb{C}\), both factors split:
\[
t^2+1=(t-i)(t+i),
\qquad
t^2+4=(t-2i)(t+2i).
\]
So
\[
\mathbb{C}[t]/\bigl((t^2+1)(t^2+4)\bigr)
\cong
\mathbb{C}\times \mathbb{C}\times \mathbb{C}\times \mathbb{C}.
\]

Thus after complexification, the two real closed points with residue field \(\mathbb{C}\) split into four complex points.

\subsubsection{Tensoring and linear factors}

By contrast, if
\[
A=\mathbb{R}[t]/\bigl((t-1)(t-2)\bigr)\cong \mathbb{R}\times \mathbb{R},
\]
then
\[
A\otimes_{\mathbb{R}}\mathbb{C}
\cong
(\mathbb{R}\times \mathbb{R})\otimes_{\mathbb{R}}\mathbb{C}
\cong
\mathbb{C}\times \mathbb{C}.
\]

Indeed,
\[
\bigl(\mathbb{R}[t]/((t-1)(t-2))\bigr)\otimes_{\mathbb{R}}\mathbb{C}
\cong
\mathbb{C}[t]/((t-1)(t-2))
\cong
\mathbb{C}\times \mathbb{C}.
\]

So in this case the quotient already had two real components, and after tensoring with \(\mathbb{C}\) these simply become two complex components.

\subsubsection{Why tensoring helps conceptually}

Tensoring with a larger field \(K\) is often called \emph{extension of scalars} or \emph{base change}.  
It is useful because:

\begin{itemize}
	\item irreducible polynomials over \(k\) may factor over \(K\),
	\item quotient rings may decompose more visibly over \(K\),
	\item geometric objects over \(k\) may split into simpler components over \(K\).
\end{itemize}

So tensoring is a controlled way of passing to a larger coefficient field while keeping the same algebraic equations.

\subsubsection{Polynomial rings in several variables}

The same principles hold in many variables.

Let \(k\to K\) be a field extension, and let
\[
A=k[s,t].
\]
Then
\[
A\otimes_k K\cong K[s,t].
\]

If \(I\subseteq k[s,t]\) is an ideal, then
\[
\bigl(k[s,t]/I\bigr)\otimes_k K \cong K[s,t]/IK,
\]
where \(IK\) denotes the ideal generated in \(K[s,t]\) by the image of \(I\).

For example, if
\[
I=((s-a)(t-b))\subseteq k[s,t],
\]
then
\[
\bigl(k[s,t]/((s-a)(t-b))\bigr)\otimes_k K
\cong
K[s,t]/((s-a)(t-b)).
\]

So the same crossing-lines equation is simply reinterpreted over the larger field \(K\).

\subsubsection{Example: complexifying the crossing lines}

Let
\[
A=\mathbb{R}[s,t]/(st).
\]
Then
\[
A\otimes_{\mathbb{R}}\mathbb{C}
\cong
\mathbb{C}[s,t]/(st).
\]

Nothing new factors here, because \(st\) was already split into linear factors over \(\mathbb{R}\).  
So base change to \(\mathbb{C}\) does not simplify the equation further.

This contrasts with irreducible quadratic factors, which do split after complexification.

\subsubsection{Tensoring and residue fields}

If \(x\in \Spec(A)\) corresponds to a prime ideal \(\mathfrak p\), its residue field is
\[
k(x)=\Frac(A/\mathfrak p)
\]
in the affine domain case, or more generally
\[
k(x)=A_{\mathfrak p}/\mathfrak p A_{\mathfrak p}.
\]

When one tensors with a field extension \(K\), the residue field may enlarge or split.

For example, over \(\mathbb{R}\), the point defined by
\[
(t^2+1)
\]
has residue field
\[
\mathbb{R}[t]/(t^2+1)\cong \mathbb{C}.
\]
After base change to \(\mathbb{C}\), this point splits into two \(\mathbb{C}\)-points corresponding to
\[
(t-i)
\qquad\text{and}\qquad
(t+i).
\]

So tensoring explains geometrically why a non-\(\mathbb{R}\)-rational point may become several \(K\)-rational points after extending scalars.

\subsubsection{A key algebraic formula}

It is useful to record the central formula again:

\[
(A/I)\otimes_A B \cong B/IB.
\]

This may be specialized in several important ways.

\paragraph{Polynomial case.}
If \(A=k[x_1,\dots,x_n]\) and \(B=K[x_1,\dots,x_n]\), then
\[
\bigl(k[x_1,\dots,x_n]/I\bigr)\otimes_k K
\cong
K[x_1,\dots,x_n]/IK.
\]

\paragraph{One polynomial.}
If \(f\in k[t]\), then
\[
\bigl(k[t]/(f)\bigr)\otimes_k K
\cong
K[t]/(f).
\]

\paragraph{Residue field extension.}
If \(k(x)\) is the residue field of a point, then tensoring \(k(x)\) with a larger base field can reveal whether that point splits over the larger field.

\subsubsection{How ideals behave under tensoring}

Let \(I=(f_1,\dots,f_r)\subseteq k[x_1,\dots,x_n]\).  
After extending scalars to \(K\), the ideal becomes
\[
IK=(f_1,\dots,f_r)\subseteq K[x_1,\dots,x_n].
\]

So tensoring does not change the formal equations; it changes only the coefficient field in which one interprets them.

However, the behavior of those equations may change dramatically, because the larger field may contain new roots or allow new factorizations.

For example:

\[
(t^2+1)\subseteq \mathbb{R}[t]
\]
is prime, because \(t^2+1\) is irreducible over \(\mathbb{R}\).

But after extending scalars to \(\mathbb{C}\), we get
\[
(t^2+1)\subseteq \mathbb{C}[t],
\]
and now
\[
t^2+1=(t-i)(t+i),
\]
so the ideal is no longer prime.

Thus tensoring can turn a prime ideal into a reducible ideal after base change.

\subsubsection{Geometric meaning of tensoring}

From the geometric point of view, tensoring corresponds to changing the base field.

If \(X=\Spec(A)\) is a \(k\)-scheme and \(K/k\) is a field extension, then the base change of \(X\) to \(K\) is
\[
X_K = X\times_{\Spec(k)} \Spec(K),
\]
and on coordinate rings this corresponds to
\[
A\otimes_k K.
\]

So all the algebraic identities above are really descriptions of geometric base change.

For example:

\begin{itemize}
	\item
	\[
	\Spec(\mathbb{R}[t]/(t^2+1))
	\]
	is one closed point over \(\mathbb{R}\) with residue field \(\mathbb{C}\),
	
	\item after base change to \(\mathbb{C}\), we get
	\[
	\Spec(\mathbb{C}[t]/(t^2+1))
	\cong
	\Spec(\mathbb{C}\times \mathbb{C}),
	\]
	which is two \(\mathbb{C}\)-points.
\end{itemize}

This is one of the clearest examples of why tensor products are central in algebraic geometry.

\subsubsection{A compact summary of the main examples}

\paragraph{Example 1.}
\[
\mathbb{R}[t]/(t-a)\cong \mathbb{R}.
\]
Tensoring with \(\mathbb{C}\) gives
\[
\bigl(\mathbb{R}[t]/(t-a)\bigr)\otimes_{\mathbb{R}}\mathbb{C}
\cong
\mathbb{C}[t]/(t-a)
\cong
\mathbb{C}.
\]

\paragraph{Example 2.}
\[
\mathbb{R}[t]/(t^2+1)\cong \mathbb{C}.
\]
Tensoring with \(\mathbb{C}\) gives
\[
\mathbb{C}\otimes_{\mathbb{R}}\mathbb{C}
\cong
\mathbb{C}[t]/(t^2+1)
\cong
\mathbb{C}\times \mathbb{C}.
\]

\paragraph{Example 3.}
\[
\mathbb{R}[t]/((t-1)(t-2))\cong \mathbb{R}\times \mathbb{R}.
\]
Tensoring with \(\mathbb{C}\) gives
\[
(\mathbb{R}\times\mathbb{R})\otimes_{\mathbb{R}}\mathbb{C}
\cong
\mathbb{C}\times \mathbb{C}.
\]

\paragraph{Example 4.}
\[
\mathbb{R}[s,t]/(st)
\]
describes two crossing real lines, and
\[
\bigl(\mathbb{R}[s,t]/(st)\bigr)\otimes_{\mathbb{R}}\mathbb{C}
\cong
\mathbb{C}[s,t]/(st).
\]

\paragraph{Example 5.}
\[
\mathbb{R}[t]/\bigl((t^2+1)(t^2+4)\bigr)\cong \mathbb{C}\times \mathbb{C},
\]
and after tensoring with \(\mathbb{C}\),
\[
\mathbb{C}[t]/\bigl((t^2+1)(t^2+4)\bigr)\cong \mathbb{C}^4.
\]

\subsubsection{Final conclusion}

Polynomial rings, ideals, quotients, and tensor products fit together through a single powerful principle:
\[
(A/I)\otimes_A B \cong B/IB.
\]

For polynomial rings this becomes:
\[
\bigl(k[x_1,\dots,x_n]/I\bigr)\otimes_k K
\cong
K[x_1,\dots,x_n]/IK.
\]

This explains:

\begin{itemize}
	\item why quotient rings change when the base field changes,
	\item why irreducible real polynomials may split over \(\mathbb{C}\),
	\item why residue fields may enlarge,
	\item why a point over \(\mathbb{R}\) may become several points over \(\mathbb{C}\),
	\item and why tensor products are fundamental in algebraic geometry.
\end{itemize}

\medskip

\textbf{Solution.}
To add tensoring to the theory of polynomial rings, ideals, and quotients, the key fact is that quotient and tensor product commute:
\[
(A/I)\otimes_A B \cong B/IB.
\]
Applied to polynomial rings, this gives
\[
\bigl(k[x_1,\dots,x_n]/I\bigr)\otimes_k K
\cong
K[x_1,\dots,x_n]/IK.
\]
Thus changing scalars from \(k\) to \(K\) simply replaces the coefficient field and transports the same equations. Over a larger field, irreducible polynomials may factor further, prime ideals may split, and points with residue fields larger than \(k\) may decompose into several \(K\)-rational points. This is exactly what happens with
\[
\mathbb{R}[t]/(t^2+1),
\]
which is one real closed point with residue field \(\mathbb{C}\), but after tensoring with \(\mathbb{C}\) becomes
\[
\mathbb{C}[t]/(t^2+1)\cong \mathbb{C}\times\mathbb{C}.
\]
So tensoring is the algebraic mechanism that makes base change visible.







	
	\section{Introduction}
	
	A good way to learn the algebraic language of schemes is to proceed in the following order:
	
	\begin{enumerate}[leftmargin=2em]
		\item define affine schemes,
		\item explain points as prime ideals,
		\item introduce local rings and residue fields,
		\item study examples in polynomial rings and in \(\Spec(\Z)\),
		\item explain ideals and quotient rings,
		\item compare one point with unions of lines,
		\item distinguish the base field from residue fields,
		\item then introduce \emph{base change} by tensor products.
	\end{enumerate}
	
	This order is natural because the later ideas depend on the earlier ones.  
	In particular, one should first understand what a point means scheme-theoretically before discussing why, for example,
	\[
	\R[t]/(t^2+1)\cong \C,
	\]
	or why after base change from \(\R\) to \(\C\) the same object may split into several simpler pieces.
	
	A small but important terminology remark is the following.
	
	\begin{remark}
		In this subject the correct phrase is \emph{base change}, not \emph{change of basis}.  
		Change of basis belongs to linear algebra.  
		Here we are changing the base ring or base field.
	\end{remark}
	
	\section{Affine schemes}
	
	\subsection{Definition}
	
	\begin{definition}
		Let \(A\) be a commutative ring with \(1\).  
		The \emph{affine scheme} associated to \(A\) is
		\[
		\Spec(A),
		\]
		the set of prime ideals of \(A\), equipped with the Zariski topology and a structure sheaf.
	\end{definition}
	
	For an introduction, the most important part is that a point of \(\Spec(A)\) is not just an ordinary coordinate point.  
	It is a prime ideal
	\[
	\p\subseteq A.
	\]
	
	So the fundamental dictionary is:
	
	\[
	\text{point of }\Spec(A)
	\quad\longleftrightarrow\quad
	\text{prime ideal }\p\subseteq A.
	\]
	
	\subsection{Closed sets}
	
	For an ideal \(I\subseteq A\), define
	\[
	V(I)=\{\p\in \Spec(A): I\subseteq \p\}.
	\]
	These are the closed sets in the Zariski topology.
	
	So the closed subsets of \(\Spec(A)\) are determined by ideals of \(A\).
	
	\subsection{Examples of affine schemes}
	
	Typical examples are:
	\[
	\Spec(k[t]),\qquad \Spec(k[s,t]),\qquad \Spec(\Z),
	\]
	where \(k\) is a field.
	
	\begin{itemize}[leftmargin=2em]
		\item \(\Spec(k[t])\) is the affine line over \(k\),
		\item \(\Spec(k[s,t])\) is the affine plane over \(k\),
		\item \(\Spec(\Z)\) is one of the most important examples in arithmetic geometry.
	\end{itemize}
	
	\section{Points as prime ideals}
	
	\subsection{General viewpoint}
	
	Let
	\[
	X=\Spec(A).
	\]
	A point \(x\in X\) corresponds to a prime ideal
	\[
	\p\subseteq A.
	\]
	
	This immediately shows that scheme-theoretic points are more general than ordinary geometric points.
	
	For example, in \(k[t]\):
	
	\begin{itemize}[leftmargin=2em]
		\item maximal ideals such as \((t-a)\) correspond to ordinary closed points,
		\item the zero ideal \((0)\), when \(k[t]\) is a domain, corresponds to the generic point.
	\end{itemize}
	
	\subsection{Closed points and generic points}
	
	\begin{definition}
		A point of \(\Spec(A)\) corresponding to a maximal ideal is called a \emph{closed point}.
	\end{definition}
	
	\begin{definition}
		If \(A\) is an integral domain, then the zero ideal \((0)\) is prime, and the corresponding point of \(\Spec(A)\) is the \emph{generic point} of the whole irreducible space \(\Spec(A)\).
	\end{definition}
	
	Thus in scheme theory we already see two very different kinds of points:
	
	\begin{itemize}[leftmargin=2em]
		\item closed points, corresponding to maximal ideals,
		\item generic points, corresponding to non-maximal prime ideals.
	\end{itemize}
	
	\section{Local rings and residue fields}
	
	\subsection{Local ring at a point}
	
	Let \(X=\Spec(A)\), and let \(x\in X\) correspond to the prime ideal \(\p\subseteq A\).
	
	\begin{definition}
		The \emph{local ring} of \(X\) at \(x\) is
		\[
		\OO_{X,x}=A_{\p}.
		\]
	\end{definition}
	
	This ring encodes functions defined near the point \(x\).
	
	Its unique maximal ideal is
	\[
	\m_x=\p A_{\p}.
	\]
	
	\subsection{Residue field}
	
	\begin{definition}
		The \emph{residue field} of the point \(x\) is
		\[
		k(x)=\OO_{X,x}/\m_x=A_{\p}/\p A_{\p}.
		\]
	\end{definition}
	
	This field tells us what the ``scalars at the point'' are.
	
	\subsection{The basic computational formula}
	
	\begin{proposition}
		Let \(A\) be a ring and let \(\p\subseteq A\) be a prime ideal. Then
		\[
		k(x)=A_{\p}/\p A_{\p}\cong \Frac(A/\p),
		\]
		where \(x\in \Spec(A)\) is the point corresponding to \(\p\).
	\end{proposition}
	
	\begin{proof}
		Localize the exact sequence
		\[
		0\longrightarrow \p \longrightarrow A \longrightarrow A/\p \longrightarrow 0
		\]
		at the multiplicative set \(A\setminus \p\). One gets
		\[
		A_{\p}/\p A_{\p}\cong S^{-1}(A/\p),
		\]
		where \(S\) is the image of \(A\setminus \p\) in \(A/\p\). Since \(A/\p\) is an integral domain, localizing at all nonzero elements gives its field of fractions:
		\[
		S^{-1}(A/\p)\cong \Frac(A/\p).
		\]
		Hence
		\[
		A_{\p}/\p A_{\p}\cong \Frac(A/\p).
		\]
	\end{proof}
	
	\begin{remark}
		This formula is often the fastest way to compute residue fields in concrete examples:
		\[
		\boxed{
			k(x)=\Frac(A/\p).
		}
		\]
	\end{remark}
	
	\section{First examples of residue fields}
	
	\subsection{The affine line \(\Spec(k[t])\)}
	
	Let
	\[
	X=\Spec(k[t]).
	\]
	
	\subsubsection{Closed rational point}
	
	Take
	\[
	\p=(t-a),\qquad a\in k.
	\]
	Then
	\[
	k[t]/(t-a)\cong k,
	\]
	so
	\[
	k(x)=\Frac(k[t]/(t-a))\cong k.
	\]
	
	Thus the point defined by \((t-a)\) has residue field \(k\).
	
	\subsubsection{Generic point}
	
	Take
	\[
	\p=(0).
	\]
	Then
	\[
	k(x)=\Frac(k[t])=k(t).
	\]
	
	So the generic point of the affine line has residue field the rational function field \(k(t)\).
	
	\subsection{The affine plane \(\Spec(k[s,t])\)}
	
	Let
	\[
	X=\Spec(k[s,t]).
	\]
	
	Take the maximal ideal
	\[
	\p=(s-a,t-b).
	\]
	Then
	\[
	k[s,t]/(s-a,t-b)\cong k,
	\]
	hence
	\[
	k(x)\cong k.
	\]
	
	So this is an ordinary \(k\)-rational point.
	
	\subsection{\(\Spec(\Z)\)}
	
	Let
	\[
	X=\Spec(\Z).
	\]
	
	\subsubsection{Closed point}
	
	For a prime number \(p\), take
	\[
	\p=(p).
	\]
	Then
	\[
	\Z/(p)\cong \F_p,
	\]
	so
	\[
	k(x)=\F_p.
	\]
	
	\subsubsection{Generic point}
	
	For
	\[
	\p=(0),
	\]
	we get
	\[
	k(x)=\Frac(\Z)=\Q.
	\]
	
	So in \(\Spec(\Z)\), the generic point has residue field \(\Q\), while the closed point \((p)\) has residue field \(\F_p\).
	
	\subsection{A non-rational closed point over \(\R\)}
	
	Consider
	\[
	X=\Spec(\R[t]).
	\]
	Take the prime ideal
	\[
	\p=(t^2+1).
	\]
	Since \(t^2+1\) is irreducible over \(\R\), this ideal is maximal. Thus
	\[
	k(x)=\R[t]/(t^2+1).
	\]
	But
	\[
	\R[t]/(t^2+1)\cong \C.
	\]
	So this point has residue field \(\C\), not \(\R\).
	
	This is one of the simplest and most important examples of a closed point that is not \(\R\)-rational.
	
	\section{Ideals and quotient rings}
	
	\subsection{Ideals define closed subschemes}
	
	If \(I\subseteq A\) is an ideal, then the quotient ring
	\[
	A/I
	\]
	defines the closed subscheme
	\[
	\Spec(A/I)\subseteq \Spec(A).
	\]
	
	Thus quotient rings are the algebraic mechanism for imposing equations.
	
	\subsection{One point versus one equation}
	
	It is essential to distinguish carefully between different ideals.
	
	\begin{example}[A single point]
		In \(k[s,t]\), the ideal
		\[
		(s-a,t-b)
		\]
		forces
		\[
		s=a,\qquad t=b.
		\]
		Hence
		\[
		k[s,t]/(s-a,t-b)\cong k.
		\]
		This is the coordinate ring of one point.
	\end{example}
	
	\begin{example}[A union of two lines]
		In \(k[s,t]\), the principal ideal
		\[
		\bigl((s-a)(t-b)\bigr)
		\]
		forces only
		\[
		(s-a)(t-b)=0.
		\]
		This means
		\[
		s=a
		\qquad\text{or}\qquad
		t=b,
		\]
		so the zero set is the union of the two lines
		\[
		s=a
		\qquad\text{and}\qquad
		t=b.
		\]
		The quotient ring is
		\[
		k[s,t]/\bigl((s-a)(t-b)\bigr).
		\]
		This is \emph{not} a field, because in the quotient we have
		\[
		[s-a]\neq 0,\qquad [t-b]\neq 0,\qquad [s-a][t-b]=0.
		\]
		So the quotient has nonzero zero divisors.
	\end{example}
	
	\subsection{Description of the crossing-lines quotient}
	
	The ring
	\[
	k[s,t]/\bigl((s-a)(t-b)\bigr)
	\]
	may be described as
	\[
	\left\{(u(t),v(s))\in k[t]\times k[s] : u(b)=v(a)\right\}.
	\]
	
	So an element of the quotient is a pair of polynomial functions:
	
	\begin{itemize}[leftmargin=2em]
		\item one on the line \(s=a\),
		\item one on the line \(t=b\),
	\end{itemize}
	
	which agree at the intersection point \((a,b)\).
	
	\subsection{Distinct ideals may give isomorphic quotients}
	
	Consider
	\[
	f_1=(s+t-1)(s-2t+3),
	\qquad
	f_2=(s-5)(s+t).
	\]
	Then the ideals
	\[
	(f_1)
	\qquad\text{and}\qquad
	(f_2)
	\]
	are distinct in \(k[s,t]\), because in a domain two nonzero principal ideals are equal only if their generators differ by a unit.
	
	However, both quotients
	\[
	k[s,t]/(f_1),
	\qquad
	k[s,t]/(f_2)
	\]
	are abstractly isomorphic to
	\[
	k[u,v]/(uv),
	\]
	because each describes two distinct lines meeting transversely in one point.
	
	So one must distinguish between:
	
	\begin{itemize}[leftmargin=2em]
		\item different ideals inside the same ring,
		\item quotient rings that may nevertheless be abstractly isomorphic.
	\end{itemize}
	
	\section{The base field and rational points}
	
	\subsection{Schemes over a field}
	
	When we write
	\[
	\A^n_k=\Spec(k[x_1,\dots,x_n]),
	\]
	the field \(k\) is called the \emph{base field}.
	
	A scheme over \(k\) may have points whose residue fields are larger than \(k\).
	
	\subsection{Rational points}
	
	\begin{definition}
		Let \(X\) be a scheme over a field \(k\). A point \(x\in X\) is called a \emph{\(k\)-rational point} if
		\[
		k(x)\cong k.
		\]
	\end{definition}
	
	So in \(\Spec(k[t])\), the point \((t-a)\) with \(a\in k\) is \(k\)-rational.  
	But in \(\Spec(\R[t])\), the point \((t^2+1)\) is not \(\R\)-rational, since
	\[
	k(x)\cong \C.
	\]
	
	\subsection{Base field versus residue field}
	
	This distinction is fundamental:
	
	\begin{itemize}[leftmargin=2em]
		\item the \emph{base field} is the field over which the scheme is defined,
		\item the \emph{residue field} \(k(x)\) is attached to a particular point \(x\).
	\end{itemize}
	
	A scheme over \(\R\) may have points with residue field \(\R\), \(\C\), or more complicated extensions.
	
	\section{Why \(\C\) appears from \(\R[t]/(t^2+1)\)}
	
	Consider
	\[
	\R[t]/(t^2+1).
	\]
	In this quotient, the class \([t]\) satisfies
	\[
	[t]^2=-1.
	\]
	Since \(\R\) contains no element whose square is \(-1\), this quotient cannot be isomorphic to \(\R\).
	
	Define
	\[
	\varphi:\R[t]\longrightarrow \C,
	\qquad
	f(t)\longmapsto f(i).
	\]
	Then \(\varphi\) is surjective and
	\[
	\ker(\varphi)=(t^2+1).
	\]
	Hence
	\[
	\R[t]/(t^2+1)\cong \C.
	\]
	
	So we do not arbitrarily switch from \(\R\) to \(\C\).  
	Rather, the quotient itself naturally becomes a field isomorphic to \(\C\).
	
	\section{Base change and tensor products}
	
	\subsection{The general idea}
	
	Now we come to \emph{base change}.  
	Suppose \(k\to K\) is a field extension. If \(X\) is a scheme over \(k\), the base change of \(X\) to \(K\) is
	\[
	X_K = X\times_{\Spec(k)} \Spec(K).
	\]
	
	In the affine case, if
	\[
	X=\Spec(A)
	\]
	with \(A\) a \(k\)-algebra, then
	\[
	X_K=\Spec(A\otimes_k K).
	\]
	
	So tensor products are the algebraic mechanism for base change.
	
	\subsection{Polynomial rings under base change}
	
	If
	\[
	A=k[x_1,\dots,x_n],
	\]
	then
	\[
	A\otimes_k K \cong K[x_1,\dots,x_n].
	\]
	
	So base change simply enlarges the coefficient field while keeping the same variables.
	
	For example,
	\[
	\R[t]\otimes_{\R}\C \cong \C[t].
	\]
	
	\subsection{Quotients and tensor products}
	
	The fundamental formula is:
	
	\begin{proposition}
		Let \(A\) be a ring, \(I\subseteq A\) an ideal, and \(B\) an \(A\)-algebra. Then
		\[
		(A/I)\otimes_A B \cong B/IB.
		\]
	\end{proposition}
	
	\begin{proof}
		Define
		\[
		\Phi:(A/I)\otimes_A B \longrightarrow B/IB
		\]
		by
		\[
		(a+I)\otimes b \longmapsto ab + IB.
		\]
		This map is well defined and \(A\)-balanced, hence induces an \(A\)-linear map on the tensor product. One checks directly that it is a ring homomorphism and that it is inverse to the map
		\[
		B/IB \longrightarrow (A/I)\otimes_A B,
		\qquad
		b+IB \longmapsto (1+I)\otimes b.
		\]
		Hence
		\[
		(A/I)\otimes_A B \cong B/IB.
		\]
	\end{proof}
	
	\begin{remark}
		This formula is extremely important. It says that quotienting and extending scalars commute.
	\end{remark}
	
	\subsection{Polynomial version}
	
	If \(I\subseteq k[x_1,\dots,x_n]\), then
	\[
	\bigl(k[x_1,\dots,x_n]/I\bigr)\otimes_k K
	\cong
	K[x_1,\dots,x_n]/IK,
	\]
	where \(IK\) is the ideal generated by the same equations, now regarded inside \(K[x_1,\dots,x_n]\).
	
	\section{Examples of base change}
	
	\subsection{A linear factor over \(\R\)}
	
	Take
	\[
	A=\R[t]/(t-a).
	\]
	Then
	\[
	A\cong \R.
	\]
	After base change to \(\C\),
	\[
	A\otimes_{\R}\C
	\cong
	\bigl(\R[t]/(t-a)\bigr)\otimes_{\R}\C
	\cong
	\C[t]/(t-a)
	\cong
	\C.
	\]
	
	So a real point remains one point after complexification.
	
	\subsection{The quadratic factor \(t^2+1\)}
	
	Take
	\[
	A=\R[t]/(t^2+1).
	\]
	We know
	\[
	A\cong \C
	\]
	as an \(\R\)-algebra.
	
	Now tensor with \(\C\):
	\[
	A\otimes_{\R}\C
	\cong
	\bigl(\R[t]/(t^2+1)\bigr)\otimes_{\R}\C
	\cong
	\C[t]/(t^2+1).
	\]
	But
	\[
	t^2+1=(t-i)(t+i)
	\]
	in \(\C[t]\), so by the Chinese remainder theorem,
	\[
	\C[t]/(t^2+1)\cong \C[t]/(t-i)\times \C[t]/(t+i)\cong \C\times \C.
	\]
	
	Therefore
	\[
	\C\otimes_{\R}\C \cong \C\times \C.
	\]
	
	Geometrically, the one real closed point with residue field \(\C\) splits into two \(\C\)-rational points after base change to \(\C\).
	
	\subsection{Two real linear factors}
	
	Take
	\[
	A=\R[t]/((t-1)(t-2)).
	\]
	Then
	\[
	A\cong \R\times \R.
	\]
	After base change to \(\C\),
	\[
	A\otimes_{\R}\C
	\cong
	(\R\times \R)\otimes_{\R}\C
	\cong
	\C\times \C.
	\]
	
	So in this case the quotient already had two real components, and after complexification they simply become two complex components.
	
	\subsection{Two irreducible quadratic factors}
	
	Consider
	\[
	A=\R[t]/\bigl((t^2+1)(t^2+4)\bigr).
	\]
	Over \(\R\), the factors are comaximal, so
	\[
	A\cong \R[t]/(t^2+1)\times \R[t]/(t^2+4)\cong \C\times \C.
	\]
	
	After tensoring with \(\C\),
	\[
	A\otimes_{\R}\C
	\cong
	\C[t]/\bigl((t^2+1)(t^2+4)\bigr).
	\]
	Now
	\[
	t^2+1=(t-i)(t+i),
	\qquad
	t^2+4=(t-2i)(t+2i),
	\]
	so
	\[
	\C[t]/\bigl((t^2+1)(t^2+4)\bigr)\cong \C^4.
	\]
	
	So the two real closed points with residue field \(\C\) split into four \(\C\)-points.
	
	\subsection{Crossing lines}
	
	Let
	\[
	A=\R[s,t]/(st).
	\]
	Then
	\[
	A\otimes_{\R}\C
	\cong
	\C[s,t]/(st).
	\]
	
	Here the equation already factors over \(\R\), so nothing new splits after base change.
	
	This contrasts with irreducible quadratic factors, which do split after complexification.
	
	\section{Geometric interpretation of base change}
	
	\subsection{Changing the coefficient field}
	
	Base change from \(k\) to \(K\) means that we keep the same equations but allow coefficients and points to be interpreted over the larger field \(K\).
	
	Thus:
	\[
	X=\Spec(A)
	\quad\leadsto\quad
	X_K=\Spec(A\otimes_k K).
	\]
	
	\subsection{Splitting phenomena}
	
	Base change often makes hidden structure visible.
	
	\begin{itemize}[leftmargin=2em]
		\item an irreducible polynomial over \(k\) may factor over \(K\),
		\item a non-\(k\)-rational point may split into several \(K\)-rational points,
		\item a prime ideal over \(k\) may become reducible over \(K\).
	\end{itemize}
	
	This is exactly what happens with
	\[
	\R[t]/(t^2+1)
	\]
	when one tensors with \(\C\).
	
	\subsection{Residue fields and splitting}
	
	The residue field explains why points may split after base change.
	
	For the point defined by
	\[
	(t^2+1)\subset \R[t],
	\]
	the residue field is
	\[
	k(x)=\R[t]/(t^2+1)\cong \C.
	\]
	
	So over \(\R\), this is one closed point with residue field \(\C\).  
	After base change to \(\C\), the polynomial factors into
	\[
	(t-i)(t+i),
	\]
	and the point splits into two \(\C\)-rational points.
	
	\section{Summary of the whole picture}
	
	We have developed the theory in the following order:
	
	\begin{enumerate}[leftmargin=2em]
		\item affine schemes are spectra of rings,
		\item points are prime ideals,
		\item the local ring at a point gives the residue field,
		\item residue fields distinguish rational points, generic points, and non-rational closed points,
		\item ideals and quotient rings define closed subschemes,
		\item different ideals may describe different embeddings but isomorphic abstract quotient rings,
		\item the base field and residue fields must be carefully distinguished,
		\item base change is performed algebraically by tensor products,
		\item quotient and tensor product commute:
		\[
		(A/I)\otimes_A B \cong B/IB,
		\]
		\item after base change, irreducible polynomials may split and points may decompose into several rational points over the larger field.
	\end{enumerate}
	
	\section{Conclusion}
	
	The transition from ordinary algebra to scheme theory becomes much clearer when one follows the chain
	
	\[
	\text{rings}
	\longrightarrow
	\text{prime ideals}
	\longrightarrow
	\text{local rings}
	\longrightarrow
	\text{residue fields}
	\longrightarrow
	\text{quotients}
	\longrightarrow
	\text{base change by tensor products}.
	\]
	
	This explains in a unified way:
	
	\begin{itemize}[leftmargin=2em]
		\item why
		\[
		k[s,t]/(s-a,t-b)\cong k
		\]
		represents one point,
		
		\item why
		\[
		k[s,t]/((s-a)(t-b))
		\]
		represents two intersecting lines,
		
		\item why
		\[
		\R[t]/(t^2+1)\cong \C,
		\]
		
		\item and why after base change from \(\R\) to \(\C\), one obtains
		\[
		\C[t]/(t^2+1)\cong \C\times \C.
		\]
	\end{itemize}
	
	So tensor products are not an extra technical gadget added later.  
	They are the natural algebraic tool for understanding how schemes behave when the base field changes.
	
	\medskip
	
	\textbf{Solution.}
	A natural step-by-step exposition is:
	first define affine schemes \(\Spec(A)\), then explain points as prime ideals, then local rings and residue fields, then concrete examples in \(k[t]\), \(k[s,t]\), \(\Z\), and \(\R[t]\). After that, one should introduce ideals and quotient rings to explain how subschemes are cut out algebraically, especially the difference between a point ideal such as \((s-a,t-b)\) and a product ideal such as \(((s-a)(t-b))\). Next one distinguishes the base field from the residue field of a point, so that non-rational closed points become conceptually clear. Only then should one introduce \emph{base change}, defined affinely by tensor products:
	\[
	X=\Spec(A)\quad\mapsto\quad X_K=\Spec(A\otimes_k K).
	\]
	The key algebraic identity is
	\[
	(A/I)\otimes_A B\cong B/IB,
	\]
	which explains how equations behave under extension of scalars. This allows one to understand examples such as
	\[
	\R[t]/(t^2+1)\cong \C
	\]
	and
	\[
	\bigl(\R[t]/(t^2+1)\bigr)\otimes_{\R}\C
	\cong
	\C[t]/(t^2+1)
	\cong
	\C\times \C.
	\]
	That is the natural mathematical flow from schemes to base change.


\end{document}

